% Luc Destombe --- licence CC BY-NC-SA 4.0
% https://creativecommons.org/licenses/by-nc-sa/4.0/deed.fr
% Fichier autonome : compiler avec pdflatex, deux fois (signets, nombre de pages).
\documentclass[11pt,a4paper]{article}
\makeatletter
\newif\iflycee@unecolonne
\lycee@unecolonnetrue
\newif\iflycee@palatino
\lycee@palatinotrue

\RequirePackage[utf8]{inputenc}
\RequirePackage[T1]{fontenc}
\RequirePackage[french]{babel}
\RequirePackage{etoolbox}

\newcommand{\ShorthandsOff}{\@ifundefined{shorthandoff}{}{\shorthandoff{;:!?}}}
\newcommand{\ShorthandsOn}{\@ifundefined{shorthandon}{}{\shorthandon{;:!?}}}
\ShorthandsOff
\AtBeginDocument{\ShorthandsOn}
\AtBeginEnvironment{tikzpicture}{\ShorthandsOff}

\RequirePackage{geometry}
\RequirePackage{fancyhdr}
\RequirePackage{lastpage}
\RequirePackage{multicol}
\RequirePackage{array}
\RequirePackage{enumitem}
\RequirePackage{xcolor}
\RequirePackage{needspace}

\RequirePackage{amsmath,amssymb}
\RequirePackage{mathpazo}
\DeclareMathSizes{10.95}{10.95}{8}{5.5}
\begingroup\catcode`\_=\active \catcode`\^=\active
\gdef_#1{\sb{\scriptscriptstyle #1}}
\gdef^#1{\sp{\scriptscriptstyle\lycee@moinsepais #1}}
\endgroup
\newcommand\lycee@moins{\mathbin{%
  \setbox\z@\hbox{$\m@th\scriptscriptstyle\lycee@moinsorig$}%
  \dimen@=.55\wd\z@ \advance\dimen@-.5pt
  \hbox to.75\wd\z@{\hss$\m@th\scriptscriptstyle\vcenter{\hbox{%
    \tikz\draw[line width=.5pt,line cap=round](0,0)--(\the\dimen@,0);}}$\hss}}}
\begingroup\lccode`\~=`\- \lowercase{\endgroup\let~}\lycee@moins
\newcommand\lycee@moinsepais{\mathcode`\-="8000 }
\AtBeginDocument{\mathchardef\lycee@moinsorig=\mathcode`\-\relax}
\AtBeginDocument{\catcode`\_=12 \mathcode`\_="8000
  \catcode`\^=12 \mathcode`\^="8000 }
\newcommand\lycee@exposants{%
  \fontdimen13\textfont2=.48\fontdimen6\textfont2
  \fontdimen14\textfont2=.48\fontdimen6\textfont2
  \fontdimen15\textfont2=.48\fontdimen6\textfont2
  \fontdimen13\scriptfont2=.48\fontdimen6\scriptfont2
  \fontdimen14\scriptfont2=.48\fontdimen6\scriptfont2
  \fontdimen15\scriptfont2=.48\fontdimen6\scriptfont2}
\every@math@size\expandafter{\the\every@math@size\lycee@exposants}
\newlength{\retraitcalculs}
\setlength{\retraitcalculs}{2em}

\RequirePackage{tikz}
\usetikzlibrary{arrows.meta,calc,babel,shapes.misc}
\RequirePackage[most]{tcolorbox}
\tcbuselibrary{skins,breakable}

\definecolor{coulDef}{HTML}{1F4E79}
\definecolor{coulProp}{HTML}{7030A0}
\definecolor{coulRem}{HTML}{C55A11}
\definecolor{coulEx}{HTML}{2E7D32}
\definecolor{coulExo}{HTML}{B03A2E}
\definecolor{coulPart}{HTML}{1F4E79}
\definecolor{coulMeth}{HTML}{1B6E4F}
\definecolor{coulCourbe}{HTML}{0B5ED7}
\definecolor{coulCourbeB}{HTML}{C00000}
\definecolor{coulCourbeC}{HTML}{2E7D32}
\definecolor{coulGrille}{HTML}{8C8C8C}
\definecolor{coulRep}{HTML}{1B6E4F}
\definecolor{coulBar}{HTML}{2E7D32}

\newcommand{\R}{\mathbb{R}}
\newcommand{\Rs}{\mathbb{R}^{*}}
\renewcommand{\le}{\leqslant}
\renewcommand{\ge}{\geqslant}
\renewcommand{\approx}{\simeq}
\newcommand{\vect}[1]{\overrightarrow{#1}}
\newcommand{\norme}[1]{\left\lVert #1 \right\rVert}
\newcommand{\intervalle}[2]{\left[#1\,;#2\right]}
\RequirePackage{cancel}
\renewcommand{\CancelColor}{\color{coulExo}}
\newcommand{\fois}[1]{\times{\color{coulExo}#1}}
\newcommand{\barre}[1]{\cancel{#1}}

\tikzset{grille/.style={coulGrille,line width=0.4pt}}
\newcommand{\quadrillage}[4]{\draw[grille,step=1] (#1,#2) grid (#3,#4);}
\tikzset{
  croix/.style={cross out,draw,line width=0.6pt,inner sep=0pt,minimum size=4pt},
  croixdroite/.style={croix,rotate=45,line width=0.8pt},
}
\newcommand{\pt}[3]{\path (#1) node[croix]{} node[#2,font=\small,inner sep=2.5pt]{$#3$};}

\newcommand{\lycee@repere}[4]{%
  \draw[grille,step=1] (#1,#2) grid (#3,#4);
  \draw[-{Stealth[length=2mm]},thick] (#1,0)--({#3+0.4},0) node[below left=-1pt]{$x$};
  \draw[-{Stealth[length=2mm]},thick] (0,#2)--(0,{#4+0.4}) node[below left=-1pt]{$y$};
  \node[below left,font=\scriptsize,inner sep=1.5pt] at (0,0) {$0$};}
\newcommand{\lycee@gradx}[1]{\foreach \k/\l in {#1}{%
  \draw (\k,0.07)--(\k,-0.07) node[below,font=\scriptsize,inner sep=1.5pt]{$\l$};}}
\newcommand{\lycee@grady}[1]{\foreach \k/\l in {#1}{%
  \draw (0.07,\k)--(-0.07,\k) node[left,font=\scriptsize,inner sep=1.5pt]{$\l$};}}
\newcommand{\lycee@intff}[2]{\left[#1\,;#2\right]}
\newcommand{\lycee@intfo}[2]{\left[#1\,;#2\right[}
\newcommand{\lycee@intof}[2]{\left]#1\,;#2\right]}
\newcommand{\lycee@intoo}[2]{\left]#1\,;#2\right[}
\newcommand{\lycee@ens}[1]{\left\{#1\right\}}
\AtBeginDocument{%
  \providecommand{\repere}{\lycee@repere}%
  \providecommand{\gradx}{\lycee@gradx}%
  \providecommand{\grady}{\lycee@grady}%
  \providecommand{\Cf}{\mathcal{C}\sb{\scriptscriptstyle f}}%
  \providecommand{\Cg}{\mathcal{C}\sb{\scriptscriptstyle g}}%
  \providecommand{\intff}{\lycee@intff}%
  \providecommand{\intfo}{\lycee@intfo}%
  \providecommand{\intof}{\lycee@intof}%
  \providecommand{\intoo}{\lycee@intoo}%
  \providecommand{\ens}{\lycee@ens}}

\newlist{questions}{enumerate}{3}
\setlist[questions,1]{label=\arabic*\textdegree),leftmargin=*,itemsep=2pt,topsep=3pt}
\setlist[questions,2]{label=\alph*),leftmargin=*,itemsep=1pt,topsep=2pt}
\setlist[questions,3]{label=\roman*),leftmargin=*,itemsep=1pt,topsep=1pt}
\setlist[itemize]{label=\textbullet,leftmargin=*,itemsep=2pt,topsep=3pt}

\newcommand{\besoinplace}[1]{\par\penalty\@M\begingroup
  \dimen@=#1\relax\dimen@ii=\pagegoal\advance\dimen@ii-\pagetotal
  \ifdim\dimen@>\dimen@ii\ifdim\dimen@ii>\z@\vfil\fi\break\fi\endgroup}

\newcommand{\lycee@pied}{\thepage/\pageref{LastPage}}
\newcommand{\entete}[2]{%
  \pagestyle{fancy}\fancyhf{}%
  \fancyhead[L]{\footnotesize\color{coulPart}\textbf{#1}}%
  \fancyhead[R]{\footnotesize\color{coulPart}\textbf{#2}}%
  \fancyfoot[C]{\footnotesize\lycee@pied}%
  \renewcommand{\headrulewidth}{0.4pt}%
  \renewcommand{\headrule}{\hbox to\headwidth{%
    \color{coulPart!40}\leaders\hrule height \headrulewidth\hfill}}}

\RequirePackage{ccicons}
\newcommand{\lycee@licence}{{\scriptsize\color{gray}%
  \copyright\ Luc Destombe --- \ccbyncsa\ CC BY-NC-SA 4.0}}
\newif\iflycee@licence \lycee@licencetrue
\newcommand{\sanslicence}{\lycee@licencefalse}
\AtBeginDocument{\iflycee@licence\AddToHookNext{shipout/foreground}{%
  \put(0,0){\rlap{\kern\dimexpr1in+\hoffset+\oddsidemargin+\textwidth\relax
    \llap{\raisebox{-\dimexpr1in+\voffset+\topmargin+\headheight+\headsep
      +\textheight+\footskip\relax}{\lycee@licence}}}}}\fi}

\newcommand{\formulesserrees}{%
  \setlength{\abovedisplayskip}{3pt}\setlength{\belowdisplayskip}{3pt}%
  \setlength{\abovedisplayshortskip}{2pt}\setlength{\belowdisplayshortskip}{3pt}}

\setlength{\parindent}{0pt}

\RequirePackage[implicit=false,hidelinks,pdfencoding=auto,psdextra,bookmarksopen,
  bookmarksopenlevel=1,pdfcreator={},pdfproducer={}]{hyperref}
\RequirePackage{bookmark}
\hypersetup{pdfauthor={Luc Destombe},pdfkeywords={CC BY-NC-SA 4.0}}
\newcounter{lycee@signet}
\newcommand{\signet}[3][]{\stepcounter{lycee@signet}%
  \edef\lycee@dest{\if\relax\detokenize{#1}\relax
    signet.\arabic{lycee@signet}\else#1\fi}%
  \hypertarget{\lycee@dest}{}\bookmark[level=#2,dest=\lycee@dest]{#3}}
\pdfstringdefDisableCommands{%
  \def\R{R}\def\vect#1{#1}\def\Cf{Cf}\def\Cg{Cg}\def\fois#1{×#1}%
  \def\barre#1{#1}\def\intervalle#1#2{[#1;#2]}\def\intff#1#2{[#1;#2]}%
  \def\textsuperscript#1{#1}\def\dfrac#1#2{#1/#2}\def\frac#1#2{#1/#2}%
  \def\vec#1{#1⃗}}%
\RequirePackage[official]{eurosym}

\tcbset{exercice corrige/.style={
  enhanced,breakable,before upper=\raggedright,
  colback=white,colframe=coulExo,
  boxrule=0.8pt,arc=2pt,
  left=6pt,right=6pt,top=8pt,bottom=5pt,
  before skip=9pt,after skip=4pt,
  coltitle=white,fonttitle=\bfseries\small,
  attach boxed title to top left={xshift=8pt,yshift=-\tcboxedtitleheight/2},
  boxed title style={colback=coulExo,arc=2pt,boxrule=0pt}}}
\newcounter{exo}
\newtcolorbox{exercice}[1][]{exercice corrige,
  phantom={\signet[exercice-\arabic{exo}]{2}{Exercice \arabic{exo}}},
  title={Exercice \theexo\ifx\relax#1\relax\else{} --- #1\fi}}
\newenvironment{exo}[1][\relax]{\refstepcounter{exo}\begin{exercice}[#1]}{\end{exercice}}

\geometry{top=1.8cm,bottom=1cm,left=1cm,right=1cm,headheight=14pt,footskip=0.6cm}
\renewcommand{\lycee@pied}{\thepage}

\renewcommand{\theexo}{\ifnum\value{exo}<10 0\fi\arabic{exo}}
\tcbset{exercice corrige/.style={
  enhanced,breakable,before upper=\raggedright,
  colback=white,colframe=coulExo,
  boxrule=0.9pt,arc=2pt,
  left=8pt,right=8pt,top=8pt,bottom=6pt,
  before skip=8pt,after skip=6pt,
  coltitle=white,fonttitle=\bfseries,
  attach boxed title to top left={xshift=8pt,yshift=-\tcboxedtitleheight/2},
  boxed title style={colback=coulExo,arc=2pt,boxrule=0pt}}}

\newcommand{\entetefiche}[2]{%
  \begin{center}
    \begin{tcolorbox}[enhanced,width=0.92\linewidth,colback=white,colframe=coulPart,
        boxrule=1.4pt,arc=3pt,halign=center,top=6pt,bottom=6pt]
      {\large\bfseries\color{coulPart} CORRIGÉ --- FICHE D'EXERCICES}\\[3pt]
      {\normalsize\bfseries #1}\\[2pt]
      {\normalsize #2}
    \end{tcolorbox}
  \end{center}}

\newcounter{partie}
\newcommand{\partie}[1]{%
  \stepcounter{partie}%
  \besoinplace{8\baselineskip}\vspace{4pt}%
  \begin{tcolorbox}[enhanced,colback=coulPart!8,colframe=coulPart,boxrule=0pt,
      leftrule=3pt,arc=0pt,left=6pt,right=4pt,top=3pt,bottom=3pt,
      before skip=6pt,after skip=2pt]
    \bfseries\color{coulPart}\Roman{partie} --- #1\signet{1}{\Roman{partie} --- #1}
  \end{tcolorbox}}
\newcommand{\partiesansnum}[1]{%
  \besoinplace{8\baselineskip}\vspace{4pt}%
  \begin{tcolorbox}[enhanced,colback=coulPart!8,colframe=coulPart,boxrule=0pt,
      leftrule=3pt,arc=0pt,left=6pt,right=4pt,top=3pt,bottom=3pt,
      before skip=6pt,after skip=2pt]
    \bfseries\color{coulPart}#1\signet{1}{#1}
  \end{tcolorbox}}

\setlist[questions,1]{label=\arabic*\textdegree),leftmargin=*,itemsep=3pt,topsep=3pt}
\setlist[questions,2]{label=\alph*),leftmargin=*,itemsep=2pt,topsep=2pt}
\setlist[questions,3]{label=\roman*),leftmargin=*,itemsep=1pt,topsep=1pt}
\newlist{expr}{enumerate}{1}
\setlist[expr]{label={\Alph*\ $=$},leftmargin=*,itemsep=2pt,topsep=3pt}

\newcommand{\res}[1]{{\color{coulRep}#1}}
\newcommand{\rem}[1]{\par\smallskip{\small\itshape\color{black!60}$\rightarrow$ #1}}
\newenvironment{calculs}{\par\smallskip\noindent$\displaystyle\begin{aligned}[t]}%
  {\end{aligned}$\par\smallskip}
\newcommand{\justif}[1]{\quad\text{\small\itshape\color{black!60}#1}}
\newcommand{\cas}[1]{\par\smallskip\noindent\textbf{\color{coulExo}#1}\nobreak}
\newcommand{\explique}[2]{\par\smallskip\noindent
  \begin{minipage}[t]{0.6\linewidth}\raggedright #1\end{minipage}\hfill
  \begin{minipage}[t]{0.37\linewidth}\raggedright\leavevmode
    {\small\itshape\color{black!60}#2}\end{minipage}\par}

\setlength{\parskip}{2pt}

\makeatother
\entete{Chapitre 2 --- Intervalles et fonctions --- corrigé}{Seconde}
\setlength{\parindent}{0pt}
\setlength{\jot}{5pt}

\definecolor{coulInt}{HTML}{D62828}
\definecolor{coulIntB}{HTML}{1565C0}
\definecolor{coulIntC}{HTML}{2E7D32}
\definecolor{coulIntD}{HTML}{E67E22}
\definecolor{coulIntE}{HTML}{7B1FA2}
\newcommand{\axegrad}[2]{%
  \draw[-{Stealth[length=2mm]}] ({#1-0.5},0)--({#2+0.5},0);
  \foreach \k in {#1,...,#2} \draw (\k,0.07)--(\k,-0.07);}
\NewDocumentCommand{\morceau}{O{coulInt} O{0} m m}{%
  \draw[#1,line width=1.8pt] (#3,#2)--(#4,#2);}
\NewDocumentCommand{\crochet}{O{coulInt} O{0} m m}{%
  \ifx[#4%
    \draw[#1,line width=1.2pt] ($(#3,#2)+(2.6pt,5pt)$) -- ++(-2.6pt,0) -- ++(0,-10pt) -- ++(2.6pt,0);
  \else
    \draw[#1,line width=1.2pt] ($(#3,#2)+(-2.6pt,5pt)$) -- ++(2.6pt,0) -- ++(0,-10pt) -- ++(-2.6pt,0);
  \fi}
\NewDocumentCommand{\bornelab}{O{0} m m}{%
  \node[below=6.5pt,font=\scriptsize,inner sep=1pt] at (#2,#1) {$#3$};}
\newcommand{\surligne}[2]{\fill[yellow!45] (#1,-0.55) rectangle (#2,0.55);}

\NewDocumentCommand{\droitey}{O{black!70} m m m}{%
  \draw[#1,thick] (#2,#4)--(#3,#4) node[right,font=\scriptsize,inner sep=1.5pt]{$y=#4$};}
\NewDocumentCommand{\solx}{O{coulInt} m}{\path[#1] (#2,0) node[croix]{};}
\newenvironment{lecture}{\par\noindent\begin{minipage}[t]{0.55\linewidth}\vspace{0pt}\raggedright}%
  {\end{minipage}\hfill}
\newenvironment{figlecture}{\begin{minipage}[t]{0.43\linewidth}\vspace{0pt}\centering}%
  {\end{minipage}\par\smallskip}
\newcommand{\gris}[1]{\leavevmode{\small\itshape\color{black!60}#1}}

\newcommand{\courbeXXV}{%
  \repere{-5}{-3}{6}{4}
  \gradx{-4,-3,-2,-1,1,2,3,4,5}
  \grady{-2,-1,1,2,3}
  \draw[coulCourbe,line width=1.1pt] (-4,2)
    .. controls (-3.667,2.333) and (-3.333,3) .. (-3,3)
    .. controls (-2.667,3) and (-2.333,2.5) .. (-2,2)
    .. controls (-1.667,1.5) and (-1.333,0.5) .. (-1,0)
    .. controls (-0.667,-0.5) and (-0.333,-0.667) .. (0,-1)
    .. controls (0.333,-1.333) and (0.667,-2) .. (1,-2)
    .. controls (1.333,-2) and (1.667,-1.333) .. (2,-1)
    .. controls (2.333,-0.667) and (2.667,-0.5) .. (3,0)
    .. controls (3.333,0.5) and (3.667,1.5) .. (4,2)
    .. controls (4.333,2.5) and (4.667,2.667) .. (5,3);
  \path[coulCourbe] (-4,2) node[croix]{} (5,3) node[croix]{};
  \node[coulCourbe,font=\footnotesize] at (-4.4,3.4) {$\Cf$};}
\newcommand{\courbeXXVII}{%
  \repere{-4}{-4}{4}{5}
  \gradx{-3,-2,-1,1,2,3}
  \grady{-3,-2,-1,1,2,3,4}
  \draw[coulCourbe,line width=1.1pt] (-3,-3)
    .. controls (-2.667,-2) and (-2.333,-0.833) .. (-2,0)
    .. controls (-1.667,0.833) and (-1.333,2) .. (-1,2)
    .. controls (-0.667,2) and (-0.333,0.5) .. (0,0)
    .. controls (0.333,-0.5) and (0.667,-1) .. (1,-1)
    .. controls (1.333,-1) and (1.667,-0.833) .. (2,0)
    .. controls (2.333,0.833) and (2.667,2.667) .. (3,4);
  \path[coulCourbe] (-3,-3) node[croix]{} (3,4) node[croix]{};
  \node[coulCourbe,font=\footnotesize] at (2.1,4.3) {$\Cf$};}
\newcommand{\courbeXXXIII}{%
  \repere{-1}{-3}{11}{3}
  \gradx{1,2,3,4,5,6,7,8,9,10}
  \grady{-2,-1,1,2}
  \draw[coulCourbe,line width=1.1pt] (0,0)
    .. controls (0.333,-0.667) and (0.667,-2) .. (1,-2)
    .. controls (1.333,-2) and (1.667,-1.333) .. (2,-1)
    .. controls (2.333,-0.667) and (2.667,-0.333) .. (3,0)
    .. controls (3.333,0.333) and (3.667,0.667) .. (4,1)
    .. controls (4.333,1.333) and (4.667,2) .. (5,2)
    .. controls (5.333,2) and (5.667,1.222) .. (6,1)
    .. controls (6.667,0.556) and (7.333,0.333) .. (8,0)
    .. controls (8.667,-0.333) and (9.333,-0.667) .. (10,-1);
  \path[coulCourbe] (0,0) node[croix]{} (10,-1) node[croix]{};
  \node[coulCourbe,font=\footnotesize] at (7.2,1.2) {$\Cf$};}
\newcommand{\courbeXXXV}{%
  \repere{-5}{-3}{6}{4}
  \gradx{-4,-3,-2,-1,1,2,3,4,5}
  \grady{-2,-1,1,2,3}
  \draw[coulCourbe,line width=1.1pt] (-4,1)
    .. controls (-3.667,0.333) and (-3.333,-0.5) .. (-3,-1)
    .. controls (-2.667,-1.5) and (-2.333,-2) .. (-2,-2)
    .. controls (-1.667,-2) and (-1.333,-1.5) .. (-1,-1)
    .. controls (-0.667,-0.5) and (-0.333,0.5) .. (0,1)
    .. controls (0.333,1.5) and (0.667,1.667) .. (1,2)
    .. controls (1.333,2.333) and (1.667,3) .. (2,3)
    .. controls (2.333,3) and (2.667,2.333) .. (3,2)
    .. controls (3.333,1.667) and (3.667,1.333) .. (4,1)
    .. controls (4.333,0.667) and (4.667,0.333) .. (5,0);
  \path[coulCourbe] (-4,1) node[croix]{} (5,0) node[croix]{};
  \node[coulCourbe,font=\footnotesize] at (2,3.5) {$\mathcal{C}$};}
\newcommand{\courbeXXXVI}{%
  \repere{-4}{-2}{4}{8}
  \gradx{-3,-2,-1,1,2,3}
  \grady{-1,1,2,3,4,5,6,7}
  \draw[coulCourbe,line width=1.1pt,domain=-3:3,samples=80] plot (\x,{((\x)*(\x)*(\x)-7*(\x)+6)/2});
  \node[coulCourbe,font=\footnotesize] at (2.4,7.4) {$\Cf$};}
\newcommand{\courbesXXXVII}{%
  \repere{-4}{-2}{8}{4}
  \gradx{-3,-2,-1,1,2,3,4,5,6,7}
  \grady{-1,1,2,3}
  \draw[coulCourbe,line width=1.1pt] (-3,1)
    .. controls (-2.667,1.333) and (-2.333,1.667) .. (-2,2)
    .. controls (-1.667,2.333) and (-1.333,3) .. (-1,3)
    .. controls (-0.667,3) and (-0.333,2.333) .. (0,2)
    .. controls (0.333,1.667) and (0.667,1.333) .. (1,1)
    .. controls (1.333,0.667) and (1.667,0.333) .. (2,0)
    .. controls (2.333,-0.333) and (2.667,-1) .. (3,-1)
    .. controls (3.333,-1) and (3.667,-0.333) .. (4,0)
    .. controls (4.333,0.333) and (4.667,0.667) .. (5,1)
    .. controls (5.333,1.333) and (5.667,1.667) .. (6,2)
    .. controls (6.333,2.333) and (6.667,2.667) .. (7,3);
  \draw[coulCourbeC,line width=1.1pt] (-3,3)
    .. controls (-2.667,2.667) and (-2.333,2.333) .. (-2,2)
    .. controls (-1.667,1.667) and (-1.333,1.333) .. (-1,1)
    .. controls (-0.667,0.667) and (-0.333,0.333) .. (0,0)
    .. controls (0.333,-0.333) and (0.667,-1) .. (1,-1)
    .. controls (1.333,-1) and (1.667,-0.333) .. (2,0)
    .. controls (2.333,0.333) and (2.667,0.667) .. (3,1)
    .. controls (3.333,1.333) and (3.667,1.667) .. (4,2)
    .. controls (4.333,2.333) and (4.667,3) .. (5,3)
    .. controls (5.333,3) and (5.667,2.333) .. (6,2)
    .. controls (6.333,1.667) and (6.667,1.333) .. (7,1);
  \node[coulCourbe,font=\footnotesize,right] at (7,3) {$\Cf$};
  \node[coulCourbeC,font=\footnotesize,right] at (7,1) {$\Cg$};}
\newcommand{\courbesXXXVIII}{%
  \repere{-7}{-3}{7}{6}
  \gradx{-6,-5,-4,-3,-2,-1,1,2,3,4,5,6}
  \grady{-2,-1,1,2,3,4,5}
  \draw[coulCourbeC,line width=1.1pt] (-6,5)--(6,-1);
  \draw[coulCourbe,line width=1.1pt] (-6,1)
    .. controls (-5.667,1.667) and (-5.333,2.5) .. (-5,3)
    .. controls (-4.667,3.5) and (-4.333,3.667) .. (-4,4)
    .. controls (-3.667,4.333) and (-3.333,5) .. (-3,5)
    .. controls (-2.667,5) and (-2.333,4.222) .. (-2,4)
    .. controls (-1.333,3.556) and (-0.667,3.444) .. (0,3)
    .. controls (0.333,2.778) and (0.667,2.333) .. (1,2)
    .. controls (1.333,1.667) and (1.667,1.333) .. (2,1)
    .. controls (2.333,0.667) and (2.667,0.333) .. (3,0)
    .. controls (3.333,-0.333) and (3.667,-0.667) .. (4,-1)
    .. controls (4.333,-1.333) and (4.667,-2) .. (5,-2)
    .. controls (5.333,-2) and (5.667,-1.333) .. (6,-1);
  \node[coulCourbeC,font=\footnotesize,above right] at (-6,5) {$\Cf$};
  \node[coulCourbe,font=\footnotesize,left] at (-6,1) {$\Cg$};}

\begin{document}

\entetefiche{Chapitre 2 --- Intervalles et fonctions}{Seconde}

\partie{Droite numérique, intervalles fermés}

\begin{exo}[Représenter des intervalles fermés]
\cas{1)}\par\nopagebreak
\explique{Pour chaque intervalle, on place les deux bornes, on trace les crochets tournés
vers l'intérieur, puis on épaissit le trait entre les deux.}{intervalle fermé : les
bornes en font partie, les crochets sont tournés vers l'intervalle}

\smallskip
\begin{minipage}[t]{0.48\linewidth}\centering
\begin{tikzpicture}[x=0.5cm]
  \axegrad{-1}{6}
  \morceau{1}{4}\crochet{1}{[}\crochet{4}{]}
  \bornelab{0}{0}\bornelab{1}{1}\bornelab{4}{4}
  \node[coulInt,right,font=\scriptsize] at (6.6,0) {$I$};
\end{tikzpicture}
\end{minipage}\hfill
\begin{minipage}[t]{0.48\linewidth}\centering
\begin{tikzpicture}[x=0.5cm]
  \axegrad{-7}{1}
  \morceau{-6}{-3}\crochet{-6}{[}\crochet{-3}{]}
  \bornelab{-6}{-6}\bornelab{-3}{-3}\bornelab{0}{0}
  \node[coulInt,right,font=\scriptsize] at (1.6,0) {$J$};
\end{tikzpicture}
\end{minipage}

\medskip
\begin{minipage}[t]{0.48\linewidth}\centering
\begin{tikzpicture}[x=0.5cm]
  \axegrad{-4}{5}
  \morceau{-2}{3}\crochet{-2}{[}\crochet{3}{]}
  \bornelab{-2}{-2}\bornelab{0}{0}\bornelab{3}{3}
  \node[coulInt,right,font=\scriptsize] at (5.6,0) {$K$};
\end{tikzpicture}
\end{minipage}\hfill
\begin{minipage}[t]{0.48\linewidth}\centering
\begin{tikzpicture}[x=0.8cm]
  \axegrad{-2}{4}
  \morceau{-0.5}{2.5}\crochet{-0.5}{[}\crochet{2.5}{]}
  \bornelab{-0.5}{-\frac12}\bornelab{0}{0}\bornelab{1}{1}\bornelab{2.5}{\frac52}
  \node[coulInt,right,font=\scriptsize] at (4.6,0) {$L$};
\end{tikzpicture}
\end{minipage}

\explique{Pour $L$ : $-\dfrac12=-0{,}5$ est au milieu de $-1$ et $0$ ; $\dfrac52=2{,}5$ est
au milieu de $2$ et $3$.}{une fraction se place par son écriture décimale}
\rem{Erreur fréquente : tourner les crochets vers l'extérieur. Pour un intervalle fermé,
ils \og embrassent\fg{} le trait épais. On écrit toujours la plus petite borne à gauche :
$\intff{-6}{-3}$, et non $\intff{-3}{-6}$.}

\cas{2)}\par\nopagebreak
\explique{Le trait épais va de $-2$ à $5$, et les deux crochets sont tournés vers
l'intérieur : les deux bornes appartiennent à l'intervalle. C'est
$\res{\intff{-2}{5}}$.}{seuls $0$ et $1$ sont écrits : ils donnent l'unité, puis on
compte les graduations, deux à gauche de $0$, cinq à droite}

\cas{3)}\par\nopagebreak
\explique{a) $\res{4\in\intff{0}{4}}$ : $4$ est une borne, et le crochet est
fermé.}{crochet tourné vers l'intervalle : la borne en fait partie}
\explique{b) $\res{9\notin\intff{3}{8{,}9}}$, car $9>8{,}9$.}{$9=9{,}0$ dépasse la borne
de droite}
\explique{c) $\res{-3\in\intff{-4}{-2}}$, car $-4\le -3\le -2$.}{on compare aux deux
bornes}
\explique{d) $\res{-50\notin\intff{-49}{-5}}$, car $-50<-49$.}{$-50$ est à gauche de
$-49$ sur la droite graduée}
\rem{Avec des nombres négatifs, on pense à la droite graduée : $-50$ est plus loin de zéro
que $-49$, donc plus petit. Erreur fréquente : croire $-50>-49$ parce que $50>49$.}
\end{exo}

\begin{exo}[Appartenance]
\explique{On compare le nombre aux deux bornes ; quand le nombre est une borne, c'est le
crochet qui décide.}{méthode 2 du cours}
\explique{a) $\res{2\in\intff{-3}{7}}$, car $-3\le 2\le 7$.}{}
\explique{b) $\res{-4\notin\intff{-3}{5}}$, car $-4<-3$.}{$-4$ est à gauche de $-3$}
\explique{c) $\res{10^{-2}\in\intff{0}{1}}$, car $10^{-2}=0{,}01$ et
$0\le 0{,}01\le 1$.}{$10^{-2}=\dfrac{1}{10^{2}}=\dfrac{1}{100}$ (chapitre 1)}
\explique{d) $\res{\pi\in\intff{3{,}1}{3{,}2}}$, car $\pi\approx 3{,}14$, et
$3{,}1\le 3{,}14\le 3{,}2$.}{on passe à l'écriture décimale}
\explique{e) $\res{-3\in\intff{-6}{-3}}$ : $-3$ est la borne de droite, et le crochet est
fermé.}{crochet tourné vers l'intervalle}
\explique{f) $\res{1{,}5\notin\intff{1{,}6}{2}}$, car $1{,}5<1{,}6$.}{on compare
$1{,}50$ et $1{,}60$ : même nombre de décimales}
\end{exo}

\begin{exo}[Nombres dans un intervalle]
\cas{1)}\par\nopagebreak
\explique{$\dfrac72=3{,}5$ : l'intervalle va de $-2$ à $3{,}5$.}{on passe à l'écriture
décimale pour situer la borne}
\explique{Les entiers positifs ou nuls de cet intervalle sont $\res{0\,;1\,;2\,;3}$.}{$4$
n'y est pas, car $4>3{,}5$ ; $-2$ et $-1$ y sont, mais ils sont négatifs}
\rem{\og Positifs ou nuls\fg{} : on n'oublie pas $0$.}

\cas{2)}\par\nopagebreak
\explique{On écrit les bornes avec deux décimales : $3{,}7=3{,}70$ et $3{,}8=3{,}80$.
Entre les deux, on trouve $3{,}71$ ; $3{,}72$ ; \ldots ; $3{,}79$. Par exemple
$\res{3{,}72}$ et $\res{3{,}75}$.}{il y en a bien d'autres, comme $3{,}705$ ou $3{,}7001$ :
entre deux nombres, il y en a toujours une infinité}

\cas{3)}\par\nopagebreak
\explique{$\pi\approx 3{,}14$, donc $3\le\pi\le 4$. Les entiers $3$ et $4$ se suivent :
l'intervalle $\res{\intff{3}{4}}$ convient.}{on encadre $\pi$ par l'entier juste en
dessous et l'entier juste au-dessus}
\end{exo}

\partie{Intervalles ouverts, inégalités}

\begin{exo}[D'un intervalle à une inégalité]
\explique{Crochet tourné vers l'intérieur : la borne est incluse, symbole $\le$.
Crochet tourné vers l'extérieur : la borne est exclue, symbole $<$. Du côté de
l'infini, il n'y a pas de borne : une seule inégalité.}{propriété 1 du cours}
\explique{a) $x\in\intff{2}{7}$ s'écrit $\res{2\le x\le 7}$ : l'intervalle est
\res{fermé}.}{les deux bornes sont incluses}
\explique{b) $x\in\intoo{-1{,}5}{+\infty}$ s'écrit $\res{x>-1{,}5}$ : l'intervalle est
\res{ouvert}.}{$-1{,}5$ est exclu ; $+\infty$ n'est pas un nombre, il n'est jamais
atteint}
\explique{c) $x\in\intof{-4}{0}$ s'écrit $\res{-4<x\le 0}$ : l'intervalle est
\res{semi-ouvert}.}{$-4$ exclu, $0$ inclus : une seule borne fait partie de
l'intervalle}
\explique{d) $x\in\intoo{-\infty}{5}$ s'écrit $\res{x<5}$ : l'intervalle est
\res{ouvert}.}{$5$ est exclu}
\explique{e) $x\in\intfo{-3}{8}$ s'écrit $\res{-3\le x<8}$ : l'intervalle est
\res{semi-ouvert}.}{$-3$ inclus, $8$ exclu}
\explique{f) $x\in\intoo{0}{1}$ s'écrit $\res{0<x<1}$ : l'intervalle est
\res{ouvert}.}{aucune des deux bornes}
\rem{Avec $+\infty$ ou $-\infty$, on écrit une seule inégalité : $x>-1{,}5$, et non
$-1{,}5<x<+\infty$, car l'infini n'est pas un nombre.}
\end{exo}

\begin{exo}[D'une phrase à un intervalle]
\explique{Pour chaque phrase : l'inégalité, puis l'intervalle, puis sa représentation.
Les crochets suivent les symboles : fermé pour $\le$ et $\ge$, ouvert pour $<$, $>$ et
du côté de l'infini.}{méthode 3 du cours}

\smallskip
\renewcommand{\arraystretch}{1.3}
\begin{tabular}{@{}p{5.6cm}@{\quad}l@{\quad}>{\raggedright\arraybackslash}p{5.6cm}@{}}
1) $x\ge -1$ : \ $x\in\res{\intfo{-1}{+\infty}}$ &
\begin{tikzpicture}[x=0.42cm,baseline=-3pt]\axegrad{-3}{5}
  \morceau{-1}{5.5}\crochet{-1}{[}\bornelab{-1}{-1}\bornelab{0}{0}\end{tikzpicture} &
\gris{\og supérieurs ou égaux\fg{} : $-1$ est inclus ; pas de limite à droite}\\[16pt]
2) $-3<x<2$ : \ $x\in\res{\intoo{-3}{2}}$ &
\begin{tikzpicture}[x=0.42cm,baseline=-3pt]\axegrad{-4}{4}
  \morceau{-3}{2}\crochet{-3}{]}\crochet{2}{[}\bornelab{-3}{-3}\bornelab{0}{0}\bornelab{2}{2}\end{tikzpicture} &
\gris{\og strictement\fg{} : les deux bornes sont exclues}\\[16pt]
3) $x<4$ : \ $x\in\res{\intoo{-\infty}{4}}$ &
\begin{tikzpicture}[x=0.42cm,baseline=-3pt]\axegrad{-2}{6}
  \morceau{-2.5}{4}\crochet{4}{[}\bornelab{0}{0}\bornelab{4}{4}\end{tikzpicture} &
\gris{$4$ est exclu ; pas de limite à gauche}\\[16pt]
4) $0\le x<6$ : \ $x\in\res{\intfo{0}{6}}$ &
\begin{tikzpicture}[x=0.42cm,baseline=-3pt]\axegrad{-2}{7}
  \morceau{0}{6}\crochet{0}{[}\crochet{6}{[}\bornelab{0}{0}\bornelab{6}{6}\end{tikzpicture} &
\gris{\og positifs ou nuls\fg{} : $x\ge 0$, donc $0$ inclus ; $6$ exclu}
\end{tabular}
\rem{\og Strictement\fg{} : la borne est exclue. \og Ou égal\fg{}, \og ou nul\fg{} : la
borne est incluse. En 4), les deux conditions à la fois s'écrivent en une seule double
inégalité.}
\end{exo}

\begin{exo}[D'une inégalité à un intervalle]
\explique{On lit chaque symbole : $\le$ ou $\ge$ donne un crochet fermé, $<$ ou $>$ un
crochet ouvert ; du côté où $x$ n'est pas limité, on écrit $-\infty$ ou $+\infty$, avec
un crochet ouvert.}{la plus petite borne s'écrit toujours à gauche}

\smallskip
\renewcommand{\arraystretch}{1.3}
\begin{tabular}{@{}p{5.6cm}@{\quad}l@{\quad}>{\raggedright\arraybackslash}p{5.6cm}@{}}
a) $-4\le x\le 2$ : \ $x\in\res{\intff{-4}{2}}$ &
\begin{tikzpicture}[x=0.42cm,baseline=-3pt]\axegrad{-5}{3}
  \morceau{-4}{2}\crochet{-4}{[}\crochet{2}{]}\bornelab{-4}{-4}\bornelab{0}{0}\bornelab{2}{2}\end{tikzpicture} &
\gris{$\le$ des deux côtés : deux crochets fermés}\\[16pt]
b) $x<3$ : \ $x\in\res{\intoo{-\infty}{3}}$ &
\begin{tikzpicture}[x=0.42cm,baseline=-3pt]\axegrad{-3}{5}
  \morceau{-3.5}{3}\crochet{3}{[}\bornelab{0}{0}\bornelab{3}{3}\end{tikzpicture} &
\gris{tous les nombres à gauche de $3$, sans $3$}\\[16pt]
c) $x>5$ : \ $x\in\res{\intoo{5}{+\infty}}$ &
\begin{tikzpicture}[x=0.42cm,baseline=-3pt]\axegrad{-1}{7}
  \morceau{5}{7.5}\crochet{5}{]}\bornelab{0}{0}\bornelab{5}{5}\end{tikzpicture} &
\gris{tous les nombres à droite de $5$, sans $5$}\\[16pt]
d) $2<x<3$ : \ $x\in\res{\intoo{2}{3}}$ &
\begin{tikzpicture}[x=0.42cm,baseline=-3pt]\axegrad{-1}{5}
  \morceau{2}{3}\crochet{2}{]}\crochet{3}{[}\bornelab{0}{0}\bornelab{2}{2}\bornelab{3}{3}\end{tikzpicture} &
\gris{deux inégalités strictes : deux crochets ouverts}\\[16pt]
e) $x\ge 1{,}5$ : \ $x\in\res{\intfo{1{,}5}{+\infty}}$ &
\begin{tikzpicture}[x=0.42cm,baseline=-3pt]\axegrad{-1}{5}
  \morceau{1.5}{5.5}\crochet{1.5}{[}\bornelab{0}{0}\bornelab{1.5}{1{,}5}\end{tikzpicture} &
\gris{$1{,}5$ au milieu de $1$ et $2$ ; crochet fermé}\\[16pt]
f) $x\le 10^{-2}$ : \ $x\in\res{\intof{-\infty}{10^{-2}}}$ &
\begin{tikzpicture}[x=0.42cm,baseline=-3pt]\axegrad{-5}{2}
  \morceau{-5.5}{0.01}\crochet{0.01}{]}\bornelab{0.01}{10^{-2}}
  \node[above=6pt,font=\scriptsize,inner sep=1pt] at (0,0) {$0$};\end{tikzpicture} &
\gris{$10^{-2}=0{,}01$ : juste à droite de $0$, presque confondu avec lui à cette échelle}
\end{tabular}
\rem{Erreur fréquente en b) : écrire $\intoo{3}{-\infty}$. La plus petite borne, ici
$-\infty$, s'écrit à gauche.}
\end{exo}

\begin{exo}[Tableau à compléter]
\cas{1)}\par\nopagebreak
\explique{$\dfrac14\le x<\dfrac23$ s'écrit $x\in\res{\intfo{\frac14}{\frac23}}$.}{$\frac14$
inclus ($\le$), $\frac23$ exclu ($<$) ; on vérifie que $\frac14<\frac23$ : $0{,}25<0{,}67$}
\explique{$x>\pi$ s'écrit $x\in\res{\intoo{\pi}{+\infty}}$.}{$\pi$ est exclu ; pas de
limite à droite}
\explique{$x<-1$ s'écrit $x\in\res{\intoo{-\infty}{-1}}$.}{$-1$ est exclu ; pas de limite
à gauche}

\cas{2)}\par\nopagebreak
\explique{Sur chaque ligne, on traduit ce qui est donné : inégalité vers intervalle, ou
intervalle vers inégalité.}{crochet fermé \,$\leftrightarrow$\, $\le$ ; crochet ouvert
\,$\leftrightarrow$\, $<$}
\begin{center}
\renewcommand{\arraystretch}{1.6}
\begin{tabular}{|c|c|c|}
\hline
\textbf{Inégalité(s)} & \textbf{Intervalle} & \textbf{Représentation}\\\hline
$-1\le x\le 4$ & $\res{\intff{-1}{4}}$ &
\begin{tikzpicture}[x=0.36cm,baseline=-3pt]\axegrad{-4}{7}
  \morceau{-1}{4}\crochet{-1}{[}\crochet{4}{]}\bornelab{-1}{-1}\bornelab{0}{0}\bornelab{4}{4}\end{tikzpicture}\\\hline
$\res{1<x\le 5}$ & $\intof{1}{5}$ &
\begin{tikzpicture}[x=0.36cm,baseline=-3pt]\axegrad{-4}{7}
  \morceau{1}{5}\crochet{1}{]}\crochet{5}{]}\bornelab{0}{0}\bornelab{1}{1}\bornelab{5}{5}\end{tikzpicture}\\\hline
$-3\le x<2$ & $\res{\intfo{-3}{2}}$ &
\begin{tikzpicture}[x=0.36cm,baseline=-3pt]\axegrad{-4}{7}
  \morceau{-3}{2}\crochet{-3}{[}\crochet{2}{[}\bornelab{-3}{-3}\bornelab{0}{0}\bornelab{2}{2}\end{tikzpicture}\\\hline
$0<x<3$ & $\res{\intoo{0}{3}}$ &
\begin{tikzpicture}[x=0.36cm,baseline=-3pt]\axegrad{-4}{7}
  \morceau{0}{3}\crochet{0}{]}\crochet{3}{[}\bornelab{0}{0}\bornelab{3}{3}\end{tikzpicture}\\\hline
$x>2$ & $\res{\intoo{2}{+\infty}}$ &
\begin{tikzpicture}[x=0.36cm,baseline=-3pt]\axegrad{-4}{7}
  \morceau{2}{7.5}\crochet{2}{]}\bornelab{0}{0}\bornelab{2}{2}\end{tikzpicture}\\\hline
$x\le 1$ & $\res{\intof{-\infty}{1}}$ &
\begin{tikzpicture}[x=0.36cm,baseline=-3pt]\axegrad{-4}{7}
  \morceau{-4.5}{1}\crochet{1}{]}\bornelab{0}{0}\bornelab{1}{1}\end{tikzpicture}\\\hline
$\res{-2\le x<6}$ & $\intfo{-2}{6}$ &
\begin{tikzpicture}[x=0.36cm,baseline=-3pt]\axegrad{-4}{7}
  \morceau{-2}{6}\crochet{-2}{[}\crochet{6}{[}\bornelab{-2}{-2}\bornelab{0}{0}\bornelab{6}{6}\end{tikzpicture}\\\hline
\end{tabular}
\end{center}
\rem{Deuxième ligne : $\intof{1}{5}$ se lit \og $1$ exclu, $5$ inclus\fg{}, d'où
$1<x\le 5$. Erreur fréquente : $1\le x<5$, en lisant les crochets à l'envers.}
\end{exo}

\begin{exo}[Bornes incluses ou exclues ?]
\explique{Les crochets ne comptent que si le nombre est une borne ; sinon, on compare aux
deux bornes.}{crochet vers l'intérieur : borne incluse ; vers l'extérieur : borne
exclue}
\explique{a) $\res{6\in\intof{-2}{6}}$ : $6$ est la borne de droite, et le crochet est
tourné vers l'intérieur.}{borne incluse}
\explique{b) $\res{-2\notin\intoo{-2}{4}}$ : en $-2$, le crochet est tourné vers
l'extérieur.}{borne exclue}
\explique{c) $\res{10^{-2}\in\intoo{0}{+\infty}}$, car $10^{-2}=0{,}01$ et
$0{,}01>0$.}{l'intervalle contient tous les nombres strictement positifs}
\explique{d) $\res{\pi\in\intoo{3{,}1}{3{,}2}}$, car $3{,}1<3{,}14<3{,}2$.}{$\pi$ n'est
pas une borne : les crochets ne jouent aucun rôle}
\explique{e) $\res{-4\notin\intoo{-\infty}{-4}}$ : la borne $-4$ est exclue.}{ce sont les
nombres strictement inférieurs à $-4$}
\explique{f) $\res{-1{,}5\in\intfo{-1{,}5}{1{,}5}}$ : en $-1{,}5$, le crochet est tourné
vers l'intérieur.}{l'autre borne, $1{,}5$, est exclue : ne pas confondre les deux}
\explique{g) $\res{0\in\intoo{-\infty}{+\infty}}$ : cet intervalle est $\R$ tout
entier.}{il contient tous les nombres réels}
\explique{h) $\res{-1\notin\intof{-1}{3}}$ : en $-1$, le crochet est tourné vers
l'extérieur.}{borne exclue}
\end{exo}

\partie{Intersection, réunion, complémentaire}

\begin{exo}[Premières intersections et réunions]
\explique{Pour chaque couple, on représente $I$ (en rouge) et $J$ (en bleu) sur le même
axe ; on lit la zone commune (surlignée) pour $I\cap J$, la zone couverte pour $I\cup J$,
puis on règle chaque borne.}{méthode 4 du cours}
\cas{a)}\par\nopagebreak
\begin{center}
\begin{tikzpicture}[x=0.8cm]
  \surligne{0}{3}
  \axegrad{-3}{7}
  \morceau[coulInt][0.3]{-1.5}{3}\crochet[coulInt][0.3]{-1.5}{[}\crochet[coulInt][0.3]{3}{[}
  \morceau[coulIntB][-0.3]{0}{6}\crochet[coulIntB][-0.3]{0}{]}\crochet[coulIntB][-0.3]{6}{]}
  \bornelab[-0.3]{-1.5}{-1{,}5}\bornelab[-0.3]{0}{0}\bornelab[-0.3]{3}{3}\bornelab[-0.3]{6}{6}
  \node[coulInt,right,font=\scriptsize] at (7.6,0.3) {$I$};
  \node[coulIntB,right,font=\scriptsize] at (7.6,-0.3) {$J$};
\end{tikzpicture}
\end{center}
\explique{Les intervalles se superposent de $0$ à $3$. Or $0\notin J$ et $3\notin I$ :
ces deux nombres ne sont pas dans les deux à la fois. $I\cap J=\res{\intoo{0}{3}}$.}{une
borne de $I\cap J$ n'est incluse que si elle est dans $I$ \emph{et} dans $J$}
\explique{À eux deux, ils couvrent la droite de $-1{,}5$ (qui est dans $I$) à $6$ (qui est
dans $J$), sans trou : $I\cup J=\res{\intff{-1{,}5}{6}}$.}{$3\notin I$, mais $3\in J$ : pas
de trou en $3$}
\cas{b)}\par\nopagebreak
\begin{center}
\begin{tikzpicture}[x=0.8cm]
  \surligne{1}{3}
  \axegrad{-2}{5}
  \morceau[coulInt][0.3]{-2.5}{3}\crochet[coulInt][0.3]{3}{[}
  \morceau[coulIntB][-0.3]{1}{5.5}\crochet[coulIntB][-0.3]{1}{[}
  \bornelab[-0.3]{0}{0}\bornelab[-0.3]{1}{1}\bornelab[-0.3]{3}{3}
  \node[coulInt,right,font=\scriptsize] at (5.6,0.3) {$I$};
  \node[coulIntB,right,font=\scriptsize] at (5.6,-0.3) {$J$};
\end{tikzpicture}
\end{center}
\explique{Zone commune de $1$ à $3$ ; $1$ est dans $J$ et dans $I$, $3$ n'est pas dans
$I$ : $I\cap J=\res{\intfo{1}{3}}$.}{$1<3$, donc $1\in I$}
\explique{$I$ va jusqu'à $-\infty$, $J$ jusqu'à $+\infty$, et ils se chevauchent : ils
couvrent toute la droite. $I\cup J=\res{\R}$.}{on peut aussi écrire
$\intoo{-\infty}{+\infty}$}
\end{exo}

\begin{exo}[Attention aux bornes]
\explique{Même méthode : on représente $I$ en rouge et $J$ en bleu, puis on règle chaque
borne en se demandant si elle est dans $I$, dans $J$, dans les deux ou dans aucun.}{c'est sur les bornes que se
joue chaque question}
\cas{a)}\par\nopagebreak
\begin{center}
\begin{tikzpicture}[x=0.8cm]
  \surligne{-2}{1}
  \axegrad{-4}{3}
  \morceau[coulInt][0.3]{-2}{3.5}\crochet[coulInt][0.3]{-2}{[}
  \morceau[coulIntB][-0.3]{-3}{1}\crochet[coulIntB][-0.3]{-3}{]}\crochet[coulIntB][-0.3]{1}{]}
  \bornelab[-0.3]{-3}{-3}\bornelab[-0.3]{-2}{-2}\bornelab[-0.3]{0}{0}\bornelab[-0.3]{1}{1}
  \node[coulInt,right,font=\scriptsize] at (3.6,0.3) {$I$};
  \node[coulIntB,right,font=\scriptsize] at (3.6,-0.3) {$J$};
\end{tikzpicture}
\end{center}
\explique{Zone commune de $-2$ à $1$ ; $-2$ et $1$ sont dans $I$ et dans $J$ :
$I\cap J=\res{\intff{-2}{1}}$.}{$-3<-2\le 1$, donc $-2\in J$}
\explique{Ensemble couvert : de $-3$, exclu, jusqu'à l'infini :
$I\cup J=\res{\intoo{-3}{+\infty}}$.}{$-3$ n'est ni dans $J$ (crochet ouvert) ni dans
$I$}
\cas{b)}\par\nopagebreak
\begin{center}
\begin{tikzpicture}[x=0.8cm]
  \axegrad{-1}{7}
  \morceau[coulInt][0.3]{3}{6.5}\crochet[coulInt][0.3]{3}{[}\crochet[coulInt][0.3]{6.5}{[}
  \morceau[coulIntB][-0.3]{1}{3}\crochet[coulIntB][-0.3]{1}{]}\crochet[coulIntB][-0.3]{3}{[}
  \bornelab[-0.3]{0}{0}\bornelab[-0.3]{1}{1}\bornelab[-0.3]{3}{3}\bornelab[-0.3]{6.5}{6{,}5}
  \node[coulInt,right,font=\scriptsize] at (7.6,0.3) {$I$};
  \node[coulIntB,right,font=\scriptsize] at (7.6,-0.3) {$J$};
\end{tikzpicture}
\end{center}
\explique{Le seul nombre qui pourrait être commun est $3$ ; or $3\in I$, mais $3\notin J$.
Aucun nombre n'est dans les deux : $I\cap J=\res{\varnothing}$.}{l'ensemble vide se note
$\varnothing$}
\explique{La réunion n'a pas de trou en $3$, puisque $3\in I$ :
$I\cup J=\res{\intoo{1}{6{,}5}}$.}{$1\notin J$ et $6{,}5\notin I$ : deux bornes exclues}
\cas{c)}\par\nopagebreak
\begin{center}
\begin{tikzpicture}[x=0.8cm]
  \surligne{-2.5}{2}
  \axegrad{-2}{5}
  \morceau[coulInt][0.3]{-2.5}{2}\crochet[coulInt][0.3]{2}{]}
  \morceau[coulIntB][-0.3]{-2.5}{4}\crochet[coulIntB][-0.3]{4}{[}
  \bornelab[-0.3]{0}{0}\bornelab[-0.3]{2}{2}\bornelab[-0.3]{4}{4}
  \node[coulInt,right,font=\scriptsize] at (5.6,0.3) {$I$};
  \node[coulIntB,right,font=\scriptsize] at (5.6,-0.3) {$J$};
\end{tikzpicture}
\end{center}
\explique{Tout nombre inférieur ou égal à $2$ est strictement inférieur à $4$ : $I$ est
contenu dans $J$. Donc $I\cap J=\res{\intof{-\infty}{2}}$ et
$I\cup J=\res{\intoo{-\infty}{4}}$.}{l'intersection est le plus petit des deux, la
réunion le plus grand}
\cas{d)}\par\nopagebreak
\begin{center}
\begin{tikzpicture}[x=0.8cm]
  \axegrad{0}{8}
  \morceau[coulInt][0.3]{-0.5}{5}\crochet[coulInt][0.3]{5}{]}
  \morceau[coulIntB][-0.3]{5}{8.5}\crochet[coulIntB][-0.3]{5}{[}
  \bornelab[-0.3]{0}{0}\bornelab[-0.3]{5}{5}
  \node[coulInt,right,font=\scriptsize] at (8.6,0.3) {$I$};
  \node[coulIntB,right,font=\scriptsize] at (8.6,-0.3) {$J$};
\end{tikzpicture}
\end{center}
\explique{Les deux intervalles se touchent en $5$, qui est dans $I$ et dans $J$ : c'est le
seul nombre commun. $I\cap J=\res{\ens{5}}$.}{un seul nombre : on l'écrit entre
accolades}
\explique{Ils couvrent toute la droite : $I\cup J=\res{\R}$.}{}
\cas{e)}\par\nopagebreak
\begin{center}
\begin{tikzpicture}[x=0.8cm]
  \axegrad{-4}{3}
  \morceau[coulInt][0.3]{-1}{3.5}\crochet[coulInt][0.3]{-1}{[}
  \morceau[coulIntB][-0.3]{-4.5}{-1}\crochet[coulIntB][-0.3]{-1}{[}
  \bornelab[-0.3]{-1}{-1}\bornelab[-0.3]{0}{0}
  \node[coulInt,right,font=\scriptsize] at (3.6,0.3) {$I$};
  \node[coulIntB,right,font=\scriptsize] at (3.6,-0.3) {$J$};
\end{tikzpicture}
\end{center}
\explique{Seul $-1$ pourrait être commun ; or $-1\in I$ mais $-1\notin J$ :
$I\cap J=\res{\varnothing}$.}{comparer avec d), où les deux crochets étaient fermés}
\explique{Pas de trou en $-1$, puisque $-1\in I$ : $I\cup J=\res{\R}$.}{}
\rem{Méthode à retenir : on dessine, on lit la zone, puis on traite chaque borne une par
une. Une borne de $I\cap J$ doit être dans $I$ \emph{et} dans $J$ ; une borne de $I\cup J$
dans $I$ \emph{ou} dans $J$.}
\end{exo}

\begin{exo}[Plusieurs intervalles]
\cas{1)}\par\nopagebreak
\explique{$-2\le x\le 6$ : $I=\res{\intff{-2}{6}}$ \quad ; \quad $-4\le x<-1$ :
$J=\res{\intfo{-4}{-1}}$.}{$\le$ : crochet fermé ; $<$ : crochet ouvert}
\explique{$x\le 6$ : $L=\res{\intof{-\infty}{6}}$ \quad ; \quad $2<x$, c'est-à-dire
$x>2$ : $M=\res{\intoo{2}{+\infty}}$.}{$2<x$ se relit de droite à gauche : $x$ est plus
grand que $2$}

\cas{2)}\par\nopagebreak
\explique{On représente les cinq intervalles au-dessus d'un même axe, chacun dans sa
couleur.}{on lit ensuite chaque intersection ou réunion sur la figure}
\begin{center}
\begin{tikzpicture}[x=0.62cm]
  \axegrad{-5}{11}
  \morceau[coulInt][2.5]{-2}{6}\crochet[coulInt][2.5]{-2}{[}\crochet[coulInt][2.5]{6}{]}
  \morceau[coulIntB][2]{-4}{-1}\crochet[coulIntB][2]{-4}{[}\crochet[coulIntB][2]{-1}{[}
  \morceau[coulIntC][1.5]{1}{10}\crochet[coulIntC][1.5]{1}{]}\crochet[coulIntC][1.5]{10}{]}
  \morceau[coulIntD][1]{-5.5}{6}\crochet[coulIntD][1]{6}{]}
  \morceau[coulIntE][0.5]{2}{11.5}\crochet[coulIntE][0.5]{2}{]}
  \node[coulInt,right,font=\scriptsize] at (11.8,2.5) {$I$};
  \node[coulIntB,right,font=\scriptsize] at (11.8,2) {$J$};
  \node[coulIntC,right,font=\scriptsize] at (11.8,1.5) {$K$};
  \node[coulIntD,right,font=\scriptsize] at (11.8,1) {$L$};
  \node[coulIntE,right,font=\scriptsize] at (11.8,0.5) {$M$};
  \foreach \k in {-4,-2,-1,0,1,2,6,10} {\bornelab{\k}{\k}}
\end{tikzpicture}
\end{center}
\explique{$I\cap J$ : zone commune de $-2$ (dans $I$ et dans $J$) à $-1$ (exclu de $J$) :
$I\cap J=\res{\intfo{-2}{-1}}$.}{$-4\le -2<-1$, donc $-2\in J$}
\explique{$I\cup J$ : de $-4$ (dans $J$) à $6$ (dans $I$), sans trou, puisque les deux
intervalles se chevauchent : $I\cup J=\res{\intff{-4}{6}}$.}{}
\explique{$J\cap K$ : $J$ s'arrête avant $-1$, $K$ commence après $1$ :
$J\cap K=\res{\varnothing}$.}{aucune zone commune}
\explique{$K\cup L$ : $L$ va de $-\infty$ à $6$, $K$ continue jusqu'à $10$ (inclus) :
$K\cup L=\res{\intof{-\infty}{10}}$.}{ils se chevauchent de $1$ à $6$ : pas de trou}
\explique{$I\cap L$ : tout nombre de $I$ vérifie $x\le 6$, donc $I$ est contenu dans $L$ :
$I\cap L=I=\res{\intff{-2}{6}}$.}{l'intersection est le plus petit des deux}
\explique{$K\cup M$ : de $1$ (exclu de $K$ et de $M$) jusqu'à l'infini :
$K\cup M=\res{\intoo{1}{+\infty}}$.}{$M$ prolonge $K$ au-delà de $10$}
\end{exo}

\begin{exo}[Complémentaire]
\explique{Le complémentaire est tout le reste de la droite. Chaque crochet se retourne :
une borne de l'ensemble n'est pas dans son complémentaire, et inversement.}{méthode 5 du
cours ; en rouge l'ensemble, en bleu son complémentaire}

\smallskip
\renewcommand{\arraystretch}{1.3}
\begin{tabular}{@{}p{6.4cm}@{\quad}l@{\quad}>{\raggedright\arraybackslash}p{4.6cm}@{}}
a) $\complement_{\R} A=\res{\intof{-\infty}{3}}$ &
\begin{tikzpicture}[x=0.38cm,baseline=-3pt]\axegrad{-2}{6}
  \morceau[coulInt][0.3]{3}{6.5}\crochet[coulInt][0.3]{3}{]}
  \morceau[coulIntB][-0.3]{-2.5}{3}\crochet[coulIntB][-0.3]{3}{]}
  \bornelab[-0.3]{0}{0}\bornelab[-0.3]{3}{3}\end{tikzpicture} &
\gris{$3\notin A$, donc $3$ est dans le complémentaire}\\[22pt]
b) $\complement_{\R} B=\res{\intoo{4}{+\infty}}$ &
\begin{tikzpicture}[x=0.38cm,baseline=-3pt]\axegrad{-2}{7}
  \morceau[coulInt][0.3]{-2.5}{4}\crochet[coulInt][0.3]{4}{]}
  \morceau[coulIntB][-0.3]{4}{7.5}\crochet[coulIntB][-0.3]{4}{]}
  \bornelab[-0.3]{0}{0}\bornelab[-0.3]{4}{4}\end{tikzpicture} &
\gris{$4\in B$, donc $4$ n'est pas dans le complémentaire}\\[22pt]
c) $\complement_{\R} C=\res{\intoo{-\infty}{-2}}$ &
\begin{tikzpicture}[x=0.38cm,baseline=-3pt]\axegrad{-5}{3}
  \morceau[coulInt][0.3]{-2}{3.5}\crochet[coulInt][0.3]{-2}{[}
  \morceau[coulIntB][-0.3]{-5.5}{-2}\crochet[coulIntB][-0.3]{-2}{[}
  \bornelab[-0.3]{-2}{-2}\bornelab[-0.3]{0}{0}\end{tikzpicture} &
\gris{$-2\in C$ : exclu du complémentaire}\\[22pt]
d) $\complement_{\R} D=\res{\intfo{6}{+\infty}}$ &
\begin{tikzpicture}[x=0.38cm,baseline=-3pt]\axegrad{-1}{8}
  \morceau[coulInt][0.3]{-1.5}{6}\crochet[coulInt][0.3]{6}{[}
  \morceau[coulIntB][-0.3]{6}{8.5}\crochet[coulIntB][-0.3]{6}{[}
  \bornelab[-0.3]{0}{0}\bornelab[-0.3]{6}{6}\end{tikzpicture} &
\gris{$6\notin D$ : inclus dans le complémentaire}\\[22pt]
e) $\complement_{\R} E=\res{\intoo{-\infty}{1}\cup\intoo{2}{+\infty}}$ &
\begin{tikzpicture}[x=0.38cm,baseline=-3pt]\axegrad{-3}{5}
  \morceau[coulInt][0.3]{1}{2}\crochet[coulInt][0.3]{1}{[}\crochet[coulInt][0.3]{2}{]}
  \morceau[coulIntB][-0.3]{-3.5}{1}\crochet[coulIntB][-0.3]{1}{[}
  \morceau[coulIntB][-0.3]{2}{5.5}\crochet[coulIntB][-0.3]{2}{]}
  \bornelab[-0.3]{0}{0}\bornelab[-0.3]{1}{1}\bornelab[-0.3]{2}{2}\end{tikzpicture} &
\gris{deux morceaux, de part et d'autre de $E$ ; $1$ et $2$ sont dans $E$}\\[22pt]
f) $\complement_{\R} F=\res{\intof{-\infty}{-5}\cup\intfo{8}{+\infty}}$ &
\begin{tikzpicture}[x=0.3cm,baseline=-3pt]\axegrad{-7}{10}
  \morceau[coulInt][0.3]{-5}{8}\crochet[coulInt][0.3]{-5}{]}\crochet[coulInt][0.3]{8}{[}
  \morceau[coulIntB][-0.3]{-7.5}{-5}\crochet[coulIntB][-0.3]{-5}{]}
  \morceau[coulIntB][-0.3]{8}{10.5}\crochet[coulIntB][-0.3]{8}{[}
  \bornelab[-0.3]{-5}{-5}\bornelab[-0.3]{0}{0}\bornelab[-0.3]{8}{8}\end{tikzpicture} &
\gris{$-5$ et $8$ ne sont pas dans $F$ : ils sont dans le complémentaire}
\end{tabular}
\rem{Erreur fréquente en e) et f) : oublier l'un des deux morceaux. Le complémentaire d'un
intervalle borné des deux côtés est toujours fait de deux morceaux.}
\end{exo}

\begin{exo}[Inégalités et intervalles : « et », « ou »]
\explique{\og ou\fg{} : au moins l'une des deux conditions, c'est une réunion ($\cup$).
\og et\fg{} : les deux conditions à la fois, c'est une intersection ($\cap$).}{on
traduit chaque condition séparément, puis on les relie}
\cas{a)}\par\nopagebreak
\explique{$-3<x\le -1$ donne $\intof{-3}{-1}$ ; $2\le x<5$ donne $\intfo{2}{5}$. Avec
\og ou\fg{} : $\res{\intof{-3}{-1}\cup\intfo{2}{5}}$.}{il y a un trou entre $-1$ et $2$ :
on ne peut pas l'écrire avec un seul intervalle}
\cas{b)}\par\nopagebreak
\explique{$\intff{-2}{0}$ donne $-2\le x\le 0$ ; $\intfo{3}{+\infty}$ donne $x\ge 3$. La
réunion se lit \og ou\fg{} : $\res{-2\le x\le 0 \text{ ou } x\ge 3}$.}{$\cup$ se traduit
par \og ou\fg}
\cas{c)}\par\nopagebreak
\explique{$x<-1$ donne $\intoo{-\infty}{-1}$ ; $x>4$ donne $\intoo{4}{+\infty}$ :
$\res{\intoo{-\infty}{-1}\cup\intoo{4}{+\infty}}$.}{c'est le complémentaire de
$\intff{-1}{4}$}
\cas{d)}\par\nopagebreak
\explique{$\intof{-\infty}{-3}$ donne $x\le -3$ ; $\intoo{2}{+\infty}$ donne $x>2$ :
$\res{x\le -3 \text{ ou } x>2}$.}{}
\cas{e)}\par\nopagebreak
\explique{$x>-2$ donne $\intoo{-2}{+\infty}$ ; $x\le 3$ donne $\intof{-\infty}{3}$. Avec
\og et\fg{} : $\intoo{-2}{+\infty}\cap\intof{-\infty}{3}=\res{\intof{-2}{3}}$.}{les deux
conditions à la fois s'écrivent $-2<x\le 3$}
\cas{f)}\par\nopagebreak
\explique{$\intfo{1}{+\infty}$ donne $x\ge 1$ ; $\intoo{-\infty}{5}$ donne $x<5$ :
$\res{x\ge 1 \text{ et } x<5}$, soit plus simplement $1\le x<5$, c'est-à-dire
$\res{\intfo{1}{5}}$.}{$\cap$ se traduit par \og et\fg}
\cas{g)}\par\nopagebreak
\explique{$\intfo{-4}{+\infty}\cap\intfo{1}{+\infty}=\res{\intfo{1}{+\infty}}$ ; plus
simplement : $\res{x\ge 1}$.}{un nombre supérieur ou égal à $1$ est forcément supérieur
à $-4$ : la condition $x\ge -4$ n'ajoute rien}
\cas{h)}\par\nopagebreak
\explique{$\intoo{-\infty}{0}\cup\intof{-\infty}{2}=\res{\intof{-\infty}{2}}$ ; plus
simplement : $\res{x\le 2}$.}{un nombre strictement négatif est aussi inférieur à $2$ :
le premier intervalle est contenu dans le second}
\cas{i)}\par\nopagebreak
\explique{$\intoo{-\infty}{-1}\cap\intoo{3}{+\infty}=\res{\varnothing}$.}{aucun nombre
n'est à la fois inférieur à $-1$ et supérieur à $3$}

\smallskip
\begin{minipage}[t]{0.32\linewidth}\centering
g)\quad\begin{tikzpicture}[x=0.32cm,baseline=-3pt]
  \surligne{1}{4.5}
  \axegrad{-5}{4}
  \morceau[coulInt][0.3]{-4}{4.5}\crochet[coulInt][0.3]{-4}{[}
  \morceau[coulIntB][-0.3]{1}{4.5}\crochet[coulIntB][-0.3]{1}{[}
  \bornelab[-0.3]{-4}{-4}\bornelab[-0.3]{0}{0}\bornelab[-0.3]{1}{1}
\end{tikzpicture}
\end{minipage}\hfill
\begin{minipage}[t]{0.32\linewidth}\centering
h)\quad\begin{tikzpicture}[x=0.32cm,baseline=-3pt]
  \axegrad{-5}{4}
  \morceau[coulInt][0.3]{-5.5}{0}\crochet[coulInt][0.3]{0}{[}
  \morceau[coulIntB][-0.3]{-5.5}{2}\crochet[coulIntB][-0.3]{2}{]}
  \bornelab[-0.3]{0}{0}\bornelab[-0.3]{2}{2}
\end{tikzpicture}
\end{minipage}\hfill
\begin{minipage}[t]{0.32\linewidth}\centering
i)\quad\begin{tikzpicture}[x=0.32cm,baseline=-3pt]
  \axegrad{-5}{4}
  \morceau[coulInt][0.3]{-5.5}{-1}\crochet[coulInt][0.3]{-1}{[}
  \morceau[coulIntB][-0.3]{3}{4.5}\crochet[coulIntB][-0.3]{3}{]}
  \bornelab[-0.3]{-1}{-1}\bornelab[-0.3]{0}{0}\bornelab[-0.3]{3}{3}
\end{tikzpicture}
\end{minipage}

\medskip
Le tableau complété :
\begin{center}
\renewcommand{\arraystretch}{1.45}
\small
\begin{tabular}{|c|c|c|c|}
\hline
 & \textbf{Inégalités} & \textbf{Intervalles} & \textbf{Plus simplement}\\\hline
a) & $-3<x\le -1$ \ ou \ $2\le x<5$ & $\intof{-3}{-1}\cup\intfo{2}{5}$ & \\\hline
b) & $-2\le x\le 0$ \ ou \ $x\ge 3$ & $\intff{-2}{0}\cup\intfo{3}{+\infty}$ & \\\hline
c) & $x<-1$ \ ou \ $x>4$ & $\intoo{-\infty}{-1}\cup\intoo{4}{+\infty}$ & \\\hline
d) & $x\le -3$ \ ou \ $x>2$ & $\intof{-\infty}{-3}\cup\intoo{2}{+\infty}$ & \\\hline
e) & $x>-2$ \ et \ $x\le 3$ & $\intoo{-2}{+\infty}\cap\intof{-\infty}{3}$ & $-2<x\le 3$ ; \ $\intof{-2}{3}$\\\hline
f) & $x\ge 1$ \ et \ $x<5$ & $\intfo{1}{+\infty}\cap\intoo{-\infty}{5}$ & $1\le x<5$ ; \ $\intfo{1}{5}$\\\hline
g) & $x\ge -4$ \ et \ $x\ge 1$ & $\intfo{-4}{+\infty}\cap\intfo{1}{+\infty}$ & $x\ge 1$ ; \ $\intfo{1}{+\infty}$\\\hline
h) & $x<0$ \ ou \ $x\le 2$ & $\intoo{-\infty}{0}\cup\intof{-\infty}{2}$ & $x\le 2$ ; \ $\intof{-\infty}{2}$\\\hline
i) & $x<-1$ \ et \ $x>3$ & $\intoo{-\infty}{-1}\cap\intoo{3}{+\infty}$ & $\varnothing$\\\hline
\end{tabular}
\end{center}
\rem{Erreur fréquente en h) : répondre $\intoo{-\infty}{0}$. Avec \og ou\fg{}, il suffit
d'une des deux conditions : $1$ vérifie $x\le 2$, il fait donc partie de l'ensemble.}
\end{exo}

\partie{Image et antécédent}

\begin{exo}[Calculer des images]
\cas{1)}\par\nopagebreak
\explique{On remplace $x$ par le nombre, entre parenthèses s'il est négatif, puis on
calcule en respectant les priorités : la puissance, puis les produits, puis la
somme.}{méthode 6 du cours}
\begin{calculs}
f(1) &= 2\times 1-3\times 1^{2} &&\justif{on remplace $x$ par $1$}\\
f(1) &= 2-3 &&\justif{$1^{2}=1$}\\
f(1) &= \res{-1} &&
\end{calculs}
\begin{calculs}
f(-2) &= 2\times(-2)-3\times(-2)^{2} &&\justif{$-2$ entre parenthèses}\\
f(-2) &= -4-3\times 4 &&\justif{la puissance d'abord : $(-2)^{2}=4$}\\
f(-2) &= -4-12 &&\justif{puis le produit : $3\times 4=12$}\\
f(-2) &= \res{-16} &&\justif{deux nombres négatifs : on ajoute $4$ et $12$}
\end{calculs}
\begin{calculs}
f\!\left(\tfrac12\right) &= 2\times\dfrac12-3\times\left(\dfrac12\right)^{2} &&\justif{on remplace $x$ par $\frac12$}\\
f\!\left(\tfrac12\right) &= 1-3\times\dfrac14 &&\justif{$2\times\frac12=1$ et $\left(\frac12\right)^{2}=\frac{1\times 1}{2\times 2}=\frac14$}\\
f\!\left(\tfrac12\right) &= \dfrac{1\fois{4}}{1\fois{4}}-\dfrac34 &&\justif{$3\times\frac14=\frac34$ ; on écrit $1$ en quarts}\\
f\!\left(\tfrac12\right) &= \res{\dfrac14} &&\justif{$\frac{4-3}{4}$}
\end{calculs}
\rem{Erreur fréquente : $-3\times(-2)^{2}=+12$. Le carré se calcule d'abord,
$(-2)^{2}=4$, puis on multiplie par $-3$ : on obtient $-12$.}

\cas{2)}\par\nopagebreak
\explique{$2$ est un antécédent de $-8$ si l'image de $2$ vaut $-8$ : on calcule
$f(2)$.}{méthode 7 : vérifier un antécédent, c'est calculer une image}
\begin{calculs}
f(2) &= 2\times 2-3\times 2^{2} &&\justif{on remplace $x$ par $2$}\\
f(2) &= 4-12 &&\justif{$2^{2}=4$ et $3\times 4=12$}\\
f(2) &= -8 &&
\end{calculs}
\explique{$f(2)=-8$ : \res{oui}, $2$ est un antécédent de $-8$ par $f$.}{on compare au
nombre $-8$}
\end{exo}

\begin{exo}[Image et antécédent]
\cas{1)}\par\nopagebreak
\begin{calculs}
g(-3) &= 4\times(-3)-5 &&\justif{on remplace $x$ par $-3$, entre parenthèses}\\
g(-3) &= -12-5 &&\justif{signes contraires : le produit est négatif}\\
g(-3) &= \res{-17} &&
\end{calculs}
\begin{calculs}
g(5) &= 4\times 5-5 &&\justif{on remplace $x$ par $5$}\\
g(5) &= \res{15} &&\justif{$20-5$}
\end{calculs}
\begin{calculs}
g\!\left(\tfrac34\right) &= 4\times\dfrac34-5 &&\justif{on remplace $x$ par $\frac34$}\\
g\!\left(\tfrac34\right) &= 3-5 &&\justif{$4\times\frac34=\frac{\barre{4}\times 3}{\barre{4}}=3$}\\
g\!\left(\tfrac34\right) &= \res{-2} &&
\end{calculs}

\cas{2)}\par\nopagebreak
\explique{On calcule l'image de $\frac12$ et on la compare à $3$.}{vérifier un antécédent
: calculer une image}
\begin{calculs}
g\!\left(\tfrac12\right) &= 4\times\dfrac12-5 &&\justif{on remplace $x$ par $\frac12$}\\
g\!\left(\tfrac12\right) &= 2-5 &&\\
g\!\left(\tfrac12\right) &= -3 &&
\end{calculs}
\explique{$g\!\left(\frac12\right)=-3$, et $-3\ne 3$ : \res{non}, $\frac12$ n'est pas un
antécédent de $3$.}{on compare au nombre $3$}

\cas{3)}\par\nopagebreak
\explique{Chercher l'antécédent de $3$, c'est résoudre l'équation $g(x)=3$.}{méthode 7 :
chercher un antécédent, c'est résoudre une équation}
\begin{calculs}
4x-5 &= 3 &&\justif{on résout $g(x)=3$}\\
4x &= 8 &&\justif{on ajoute $5$ aux deux membres}\\
x &= 2 &&\justif{on divise par $4$}
\end{calculs}
\explique{L'antécédent de $3$ par $g$ est $\res{2}$.}{vérification : $4\times 2-5=3$}
\rem{Question 2 : on \emph{vérifie} (on calcule une image). Question 3 : on
\emph{cherche} (on résout une équation). Les deux démarches ne se confondent pas.}
\end{exo}

\begin{exo}[Calculer des antécédents]
\explique{Dans chaque cas, on résout l'équation $f(x)=y$.}{chercher un antécédent, c'est
résoudre une équation}
\cas{a)}\par\nopagebreak
\begin{calculs}
x-4 &= -2 &&\justif{on résout $f(x)=-2$}\\
x &= 2 &&\justif{on ajoute $4$ aux deux membres}
\end{calculs}
\explique{L'antécédent de $-2$ est $\res{2}$.}{vérification : $2-4=-2$}
\cas{b)}\par\nopagebreak
\explique{On résout $x^{2}=36$ : on cherche les nombres dont le carré vaut $36$. Or
$6^{2}=36$ et $(-6)^{2}=36$.}{deux nombres opposés ont le même carré}
\explique{$36$ a deux antécédents : $\res{-6}$ et $\res{6}$.}{}
\cas{c)}\par\nopagebreak
\begin{calculs}
\dfrac1x &= \dfrac41 &&\justif{on résout $f(x)=4$, avec $4=\frac41$}\\
1\times 1 &= 4x &&\justif{produit en croix}\\
x &= \dfrac14 &&\justif{on divise par $4$}
\end{calculs}
\explique{L'antécédent de $4$ est $\res{\dfrac14}$.}{vérification : $1\div\frac14=1\times 4=4$}
\cas{d)}\par\nopagebreak
\explique{On résout $x(x-3)=0$. Un produit est nul lorsque l'un de ses facteurs est nul :
$x=0$ ou $x-3=0$, c'est-à-dire $x=0$ ou $x=3$.}{rappel du collège ; on ne développe pas : la forme produit
donne directement les solutions}
\explique{$0$ a deux antécédents : $\res{0}$ et $\res{3}$.}{vérification : $0\times(-3)=0$
et $3\times 0=0$}
\rem{En b), erreur fréquente : oublier l'antécédent négatif $-6$.}
\end{exo}

\begin{exo}[Deux fonctions égales]
\cas{1)}\par\nopagebreak
\begin{calculs}
f(2) &= (2+4)(3-2)+2^{2}-10 &&\justif{on remplace $x$ par $2$}\\
f(2) &= 6\times 1+4-10 &&\justif{les parenthèses, puis la puissance}\\
f(2) &= \res{0} &&\justif{$6+4-10$}
\end{calculs}
\begin{calculs}
g(2) &= 2-2 &&\\
g(2) &= \res{0} &&
\end{calculs}
\explique{$2$ a la même image, $0$, par $f$ et par $g$.}{}

\cas{2)}\par\nopagebreak
\explique{Par exemple avec $0$, qui rend le calcul simple.}{tout nombre convient}
\begin{calculs}
f(0) &= (0+4)(3-0)+0^{2}-10 &&\justif{on remplace $x$ par $0$}\\
f(0) &= 12-10 &&\justif{$4\times 3=12$}\\
f(0) &= \res{2} &&\\
g(0) &= 2-0 &&\\
g(0) &= \res{2} &&
\end{calculs}
\explique{Encore la même image par $f$ et par $g$.}{}

\cas{3)}\par\nopagebreak
\explique{Deux exemples ne suffisent pas : pour \emph{tout} nombre $x$, on développe et
on réduit $f(x)$.}{un exemple vérifie, il ne démontre pas}
\begin{calculs}
f(x) &= (x+4)(3-x)+x^{2}-10 &&\justif{on recopie}\\
f(x) &= 3x-x^{2}+12-4x+x^{2}-10 &&\justif{chaque terme de $(x+4)$ multiplie $3$, puis $-x$}\\
f(x) &= -x^{2}+x^{2}+3x-4x+12-10 &&\justif{on regroupe les termes de même sorte}\\
f(x) &= 2-x &&\justif{$-x^{2}+x^{2}=0$ ; $3x-4x=-x$ ; $12-10=2$}
\end{calculs}
\explique{Pour tout nombre $x$, $f(x)=2-x$, c'est-à-dire $\res{f(x)=g(x)}$ : tout nombre a
la même image par $f$ et par $g$.}{les fonctions $f$ et $g$ sont égales, même si leurs
formules ne se ressemblent pas}
\rem{Dans le développement, attention aux signes : $x\times(-x)=-x^{2}$ et
$4\times(-x)=-4x$.}
\end{exo}

\partie{Ensemble de définition}

\begin{exo}[Un nombre sans image]
\explique{Pour chaque fonction, on cherche le calcul impossible : une division par zéro,
ou la racine carrée d'un nombre négatif.}{propriété 2 du cours}
\explique{$f_1$ : $f_1(0)=\dfrac10$ n'existe pas, car on ne divise pas par zéro. Le nombre
sans image est $\res{0}$.}{les autres en ont une : $f_1(2)=\frac12$, $f_1(9)=\frac19$\ldots}
\explique{$f_2$ : le dénominateur $x-2$ vaut $0$ pour $x=2$ ; $f_2(2)=\dfrac70$ n'existe
pas. Le nombre sans image est $\res{2}$.}{$f_2(-5)=\frac{0}{-7}=0$ existe : un
numérateur nul ne gêne pas}
\explique{$f_3$ : $\sqrt{-5}$ n'existe pas, car un nombre négatif n'a pas de racine carrée.
Le nombre sans image est $\res{-5}$.}{$\sqrt0=0$ et $\sqrt9=3$ existent}
\rem{Piège pour $f_2$ : $-5$ annule le numérateur, pas le dénominateur. Seul un
dénominateur nul rend le calcul impossible.}
\end{exo}

\begin{exo}[Valeurs interdites]
\explique{Pour chaque fonction, on résout l'équation \og dénominateur $=0$\fg{}.}{méthode 8
du cours}
\cas{a)}\par\nopagebreak
\begin{calculs}
x-6 &= 0 &&\justif{le dénominateur s'annule}\\
x &= 6 &&\justif{on ajoute $6$}
\end{calculs}
\explique{La valeur interdite est $\res{6}$.}{$f(6)=\frac10$ n'existe pas}
\cas{b)}\par\nopagebreak
\begin{calculs}
x+5 &= 0 &&\justif{le dénominateur s'annule}\\
x &= -5 &&\justif{on retranche $5$}
\end{calculs}
\explique{La valeur interdite est $\res{-5}$.}{$f(-5)=\frac20$ n'existe pas}
\cas{c)}\par\nopagebreak
\begin{calculs}
3x &= 0 &&\justif{le dénominateur s'annule}\\
x &= 0 &&\justif{on divise par $3$}
\end{calculs}
\explique{La valeur interdite est $\res{0}$.}{$3\times 0=0$}
\cas{d)}\par\nopagebreak
\begin{calculs}
2x-8 &= 0 &&\justif{le dénominateur s'annule}\\
2x &= 8 &&\justif{on ajoute $8$}\\
x &= 4 &&\justif{on divise par $2$}
\end{calculs}
\explique{La valeur interdite est $\res{4}$.}{vérification : $2\times 4-8=0$}
\cas{e)}\par\nopagebreak
\begin{calculs}
2x+6 &= 0 &&\justif{le dénominateur s'annule}\\
2x &= -6 &&\justif{on retranche $6$}\\
x &= -3 &&\justif{on divise par $2$}
\end{calculs}
\explique{La valeur interdite est $\res{-3}$.}{vérification : $2\times(-3)+6=0$}
\cas{f)}\par\nopagebreak
\begin{calculs}
4x-12 &= 0 &&\justif{le dénominateur s'annule}\\
4x &= 12 &&\justif{on ajoute $12$}\\
x &= 3 &&\justif{on divise par $4$}
\end{calculs}
\explique{La valeur interdite est $\res{3}$.}{vérification : $4\times 3-12=0$}
\rem{Seul le dénominateur compte. En d), le numérateur s'annule pour $x=-2$, mais
$f(-2)=\frac{0}{-12}=0$ existe : $-2$ n'est pas une valeur interdite.}
\end{exo}

\begin{exo}[Racines carrées]
\explique{Le nombre sous la racine doit être positif ou nul : on résout \og nombre sous la
racine $\ge 0$\fg{} comme une équation, puis on conclut avec un intervalle.}{méthode 9 du
cours}
\cas{a)}\par\nopagebreak
\explique{$\sqrt{x}$ existe lorsque $x\ge 0$ : $D_f=\res{\intfo{0}{+\infty}}$.}{$0$ est
inclus : $\sqrt0=0$}
\cas{b)}\par\nopagebreak
\begin{calculs}
x-2 &\ge 0 &&\justif{le nombre sous la racine est positif ou nul}\\
x &\ge 2 &&\justif{on ajoute $2$}
\end{calculs}
\explique{$D_f=\res{\intfo{2}{+\infty}}$.}{$f(2)=\sqrt0=0$ existe : crochet fermé}
\cas{c)}\par\nopagebreak
\begin{calculs}
x+4 &\ge 0 &&\justif{le nombre sous la racine est positif ou nul}\\
x &\ge -4 &&\justif{on retranche $4$}
\end{calculs}
\explique{$D_f=\res{\intfo{-4}{+\infty}}$.}{}
\cas{d)}\par\nopagebreak
\begin{calculs}
x-5 &\ge 0 &&\\
x &\ge 5 &&\justif{on ajoute $5$}
\end{calculs}
\explique{$D_f=\res{\intfo{5}{+\infty}}$.}{}
\cas{e)}\par\nopagebreak
\begin{calculs}
3x-9 &\ge 0 &&\justif{le nombre sous la racine est positif ou nul}\\
3x &\ge 9 &&\justif{on ajoute $9$}\\
x &\ge 3 &&\justif{on divise par $3$, qui est positif}
\end{calculs}
\explique{$D_f=\res{\intfo{3}{+\infty}}$.}{vérification : $3\times 3-9=0$}
\cas{f)}\par\nopagebreak
\begin{calculs}
2x+8 &\ge 0 &&\\
2x &\ge -8 &&\justif{on retranche $8$}\\
x &\ge -4 &&\justif{on divise par $2$, qui est positif}
\end{calculs}
\explique{$D_f=\res{\intfo{-4}{+\infty}}$.}{vérification : $2\times(-4)+8=0$}
\rem{La borne est toujours incluse, car $\sqrt0=0$ existe. Erreur fréquente : écrire
$\intoo{2}{+\infty}$ en b).}
\end{exo}

\begin{exo}[Pour s'entraîner]
\explique{Sans quotient ni racine, $D_f=\R$ ; avec un quotient, on retire la valeur
interdite ; avec une racine, on résout \og nombre sous la racine $\ge 0$\fg{}.}{méthode 9
du cours}
\cas{a)}\par\nopagebreak
\explique{Ni quotient ni racine : on peut calculer $3x^{2}-2$ pour tout nombre.
$D_f=\res{\R}$, c'est-à-dire $\intoo{-\infty}{+\infty}$.}{}
\cas{b)}\par\nopagebreak
\explique{$f$ est affine : $D_f=\res{\R}$.}{$5-2x=-2x+5$, de la forme $ax+b$}
\cas{c)}\par\nopagebreak
\begin{calculs}
x-1 &= 0 &&\justif{le dénominateur du quotient s'annule}\\
x &= 1 &&
\end{calculs}
\explique{On retire la valeur interdite $1$ : $D_f=\res{\intoo{-\infty}{1}\cup\intoo{1}{+\infty}}$.}{le
terme $+x$ ne pose aucun problème}
\cas{d)}\par\nopagebreak
\begin{calculs}
5x+15 &= 0 &&\justif{le dénominateur s'annule}\\
5x &= -15 &&\justif{on retranche $15$}\\
x &= -3 &&\justif{on divise par $5$}
\end{calculs}
\explique{$D_f=\res{\intoo{-\infty}{-3}\cup\intoo{-3}{+\infty}}$.}{le numérateur $2x+3$
ne compte pas}
\cas{e)}\par\nopagebreak
\begin{calculs}
4x-8 &\ge 0 &&\justif{le nombre sous la racine est positif ou nul}\\
4x &\ge 8 &&\justif{on ajoute $8$}\\
x &\ge 2 &&\justif{on divise par $4$, qui est positif}
\end{calculs}
\explique{$D_f=\res{\intfo{2}{+\infty}}$.}{le $+1$ est ajouté après la racine : il ne
change rien}
\rem{En c) et d), on retire un seul nombre : $\R$ privé d'un point s'écrit avec deux
intervalles ouverts en ce nombre.}
\end{exo}

\partie{Tableau de valeurs}

\begin{exo}[Lire un tableau]
\begin{center}
\renewcommand{\arraystretch}{1.25}
\begin{tabular}{|c|*{5}{c|}l}
\cline{1-6}
$x$ & $-2$ & $-1$ & $0$ & $1$ & $2$ & $\leftarrow$ antécédents\\\cline{1-6}
$f(x)$ & $2$ & $0$ & $-3$ & $-1$ & $4$ & $\leftarrow$ images\\\cline{1-6}
\end{tabular}
\end{center}
\cas{1)}\par\nopagebreak
\explique{Une image se lit de haut en bas : on cherche le nombre sur la ligne des $x$,
son image est juste en dessous.}{définition 10 du cours}
\explique{$f(-1)=\res{0}$ \ ; \ $f(1)=\res{-1}$ \ ; \ $f(-2)=\res{2}$.}{}
\cas{2)}\par\nopagebreak
\explique{Sous $x=2$, on lit $f(2)=4$, et non $-1$ : $2$ \res{n'est pas} un antécédent de
$-1$.}{on vérifie avec l'image de $2$}
\explique{Un antécédent se lit de bas en haut : on cherche $-1$ sur la ligne des images ;
au-dessus, on lit $1$. Le tableau donne l'antécédent $\res{1}$.}{il peut y en avoir
d'autres, que le tableau ne montre pas}
\rem{$-1$ apparaît sur les deux lignes : en haut, c'est un antécédent (son image est
$0$) ; en bas, c'est une image (son antécédent lu est $1$). On regarde la bonne ligne.}
\end{exo}

\begin{exo}[Lire un tableau (suite)]
\begin{center}
\renewcommand{\arraystretch}{1.25}
\begin{tabular}{|c|*{5}{c|}}
\hline
$x$ & $-3$ & $-1$ & $2$ & $4$ & $6$\\\hline
$g(x)$ & $4$ & $-2$ & $0$ & $5$ & $4$\\\hline
\end{tabular}
\end{center}
\cas{1)}\par\nopagebreak
\explique{$g(4)=\res{5}$ \ ; \ $g(2)=\res{0}$ \ ; \ $g(-1)=\res{-2}$.}{de haut en bas : le
nombre sur la ligne des $x$, l'image en dessous}
\cas{2)}\par\nopagebreak
\explique{Antécédent de $5$ : $\res{4}$, car $g(4)=5$.}{de bas en haut : on cherche $5$
sur la ligne des images}
\explique{Antécédent de $-2$ : $\res{-1}$, car $g(-1)=-2$.}{}
\explique{Antécédents de $4$ : $\res{-3}$ et $\res{6}$, car $g(-3)=4$ et $g(6)=4$.}{on
parcourt toute la ligne des images : $4$ y figure deux fois}
\rem{Un nombre n'a qu'une image, mais il peut avoir plusieurs antécédents. Ne pas répondre
$5$ pour les antécédents de $4$ : $5$ est l'image de $4$.}
\end{exo}

\begin{exo}[Construire un tableau]
\cas{1)}\par\nopagebreak
\explique{Premier nombre $-2$, dernier $4$, pas de $1$ : $x$ prend les valeurs $-2$ ; $-1$ ;
$0$ ; $1$ ; $2$ ; $3$ ; $4$. On calcule chaque image, la parenthèse d'abord, puis le
carré.}{le pas est l'écart entre deux valeurs de $x$ qui se suivent}
\begin{calculs}
f(-2) &= 0{,}5\times(-2-1)^{2}-2 &&\justif{on remplace $x$ par $-2$}\\
f(-2) &= 0{,}5\times 9-2 &&\justif{$-2-1=-3$, puis $(-3)^{2}=9$}\\
f(-2) &= 2{,}5 &&\justif{$4{,}5-2$}
\end{calculs}
\begin{calculs}
f(-1) &= 0{,}5\times(-2)^{2}-2 = 0 &&\justif{$(-2)^{2}=4$ et $0{,}5\times 4=2$}\\
f(0) &= 0{,}5\times(-1)^{2}-2 = -1{,}5 &&\justif{$0{,}5\times 1=0{,}5$}\\
f(1) &= 0{,}5\times 0^{2}-2 = -2 &&\justif{la parenthèse est nulle}\\
f(2) &= 0{,}5\times 1^{2}-2 = -1{,}5 &&\\
f(3) &= 0{,}5\times 2^{2}-2 = 0 &&\\
f(4) &= 0{,}5\times 3^{2}-2 = 2{,}5 &&\justif{$0{,}5\times 9=4{,}5$}
\end{calculs}
\begin{center}
\renewcommand{\arraystretch}{1.25}
\begin{tabular}{|c|*{7}{c|}}
\hline
$x$ & $-2$ & $-1$ & $0$ & $1$ & $2$ & $3$ & $4$\\\hline
$f(x)$ & $2{,}5$ & $0$ & $-1{,}5$ & $-2$ & $-1{,}5$ & $0$ & $2{,}5$\\\hline
\end{tabular}
\end{center}
\rem{$0{,}5\times(x-1)^{2}$ : on élève au carré avant de multiplier par $0{,}5$. Erreur
fréquente : $0{,}5\times(-3)=-1{,}5$, puis $(-1{,}5)^{2}$.}

\cas{2)}\par\nopagebreak
\explique{On cherche $0$ sur la ligne des images : il est sous $-1$ et sous $3$. Deux
antécédents de $0$ : $\res{-1}$ et $\res{3}$.}{de bas en haut}
\explique{De même, $2{,}5$ est sous $-2$ et sous $4$ : $\res{-2}$ et $\res{4}$.}{}

\cas{3)}\par\nopagebreak
\begin{lecture}
On place les sept points du tableau, de coordonnées $\bigl(x\,;f(x)\bigr)$, puis on les
relie à main levée par une courbe sans angle.

\smallskip
\gris{Le point le plus bas est $(1\,;-2)$ ; les images se répètent de part et d'autre de
$1$ ($f(0)=f(2)$, $f(-1)=f(3)$\ldots) : la courbe est symétrique par rapport à la droite
verticale qui passe par ce point.}
\end{lecture}
\begin{figlecture}
\begin{tikzpicture}[x=0.5cm,y=0.5cm]
  \repere{-3}{-3}{5}{4}
  \gradx{-2,-1,1,2,3,4}
  \grady{-2,-1,1,2,3}
  \draw[coulCourbe,line width=1.1pt,domain=-2:4,samples=60] plot (\x,{0.5*(\x-1)*(\x-1)-2});
  \foreach \a/\b in {-2/2.5,-1/0,0/-1.5,1/-2,2/-1.5,3/0,4/2.5} \path[coulCourbe] (\a,\b) node[croix]{};
  \node[coulCourbe,font=\footnotesize,right] at (4,2.5) {$\Cf$};
\end{tikzpicture}
\end{figlecture}
\end{exo}

\partie{Courbe représentative}

\begin{exo}[Lecture graphique]
\explique{Sur la courbe, un point a pour abscisse un antécédent $x$ et pour ordonnée son
image $f(x)$.}{définition 11 du cours}
\cas{a)}\par\nopagebreak
\begin{lecture}
La courbe commence au point d'abscisse $-4$ et finit au point d'abscisse $5$ :
$D_f=\res{\intff{-4}{5}}$.

\smallskip
\gris{On lit les abscisses des extrémités (marquées par des croix). Crochets fermés : les
extrémités sont des points de la courbe, $-4$ et $5$ ont une image.}
\end{lecture}
\begin{figlecture}
\begin{tikzpicture}[x=0.45cm,y=0.45cm]
  \courbeXXV
  \draw[black!70,dashed] (-4,2)--(-4,0) (5,3)--(5,0);
  \morceau{-4}{5}\crochet{-4}{[}\crochet{5}{]}
\end{tikzpicture}
\end{figlecture}
\cas{b) et c)}\par\nopagebreak
\begin{lecture}
b) Au-dessus de $3$, la courbe est sur l'axe des abscisses : $f(3)=\res{0}$. De même,
$f(-1)=\res{0}$.

\smallskip
\gris{Le point de la courbe d'abscisse $3$ est sur l'axe des abscisses : son ordonnée
est $0$.}

\medskip
c) On part de $-2$ sur l'axe des abscisses, on monte jusqu'à la courbe, puis on va
horizontalement jusqu'à l'axe des ordonnées : $f(-2)=\res{2}$. De même,
$f(-4)=\res{2}$.

\smallskip
\gris{Image : de l'axe des abscisses vers la courbe, puis vers l'axe des ordonnées.}
\end{lecture}
\begin{figlecture}
\begin{tikzpicture}[x=0.45cm,y=0.45cm]
  \courbeXXV
  \solx{3}\solx{-1}
  \draw[black!70,dashed] (-2,0)--(-2,2)--(0,2) (-4,0)--(-4,2)--(-2,2);
  \path[coulInt] (-2,2) node[croix]{} (-4,2) node[croix]{};
\end{tikzpicture}
\end{figlecture}
\cas{d)}\par\nopagebreak
\begin{lecture}
Les antécédents de $2$ sont les abscisses des points de la courbe situés sur la droite
horizontale $y=2$ : $\res{-4\,;-2\,;4}$.

\smallskip
\gris{Antécédent : de l'axe des ordonnées vers la courbe, puis vers l'axe des abscisses.
On n'en oublie aucun : on suit la droite d'un bout à l'autre.}

\medskip
Antécédents de $0$ : la courbe coupe l'axe des abscisses en $\res{-1}$ et $\res{3}$
(croix vertes).

Antécédent de $-2$ : la droite $y=-2$ touche la courbe en son point le plus bas, d'abscisse
$\res{1}$ (en bleu).

La droite $y=-3$ ne rencontre pas la courbe : $-3$ \res{n'a pas d'antécédent}.
\end{lecture}
\begin{figlecture}
\begin{tikzpicture}[x=0.45cm,y=0.45cm]
  \courbeXXV
  \droitey[coulInt]{-5}{6}{2}
  \draw[coulInt,dashed] (-4,2)--(-4,0) (-2,2)--(-2,0) (4,2)--(4,0);
  \solx{-4}\solx{-2}\solx{4}
  \droitey[coulIntB]{-5}{6}{-2}
  \draw[coulIntB,dashed] (1,-2)--(1,0);
  \solx[coulIntB]{1}
  \droitey[black!60]{-5}{6}{-3}
  \solx[coulCourbeC]{-1}\solx[coulCourbeC]{3}
\end{tikzpicture}
\end{figlecture}
\cas{e)}\par\nopagebreak
\begin{lecture}
$f(x)=3$ : on cherche les antécédents de $3$. La droite $y=3$ touche la courbe en son
point le plus haut à gauche, d'abscisse $-3$, et en son extrémité droite, d'abscisse $5$ :
$\res{-3}$ et $\res{5}$.

\smallskip
\gris{Vérification : les points $(-3\,;3)$ et $(5\,;3)$ sont sur la courbe.}
\end{lecture}
\begin{figlecture}
\begin{tikzpicture}[x=0.45cm,y=0.45cm]
  \courbeXXV
  \droitey[coulInt]{-5}{6}{3}
  \draw[coulInt,dashed] (-3,3)--(-3,0) (5,3)--(5,0);
  \solx{-3}\solx{5}
\end{tikzpicture}
\end{figlecture}
\rem{Erreur fréquente : lire l'image sur le mauvais axe. L'image se lit sur l'axe des
ordonnées, l'antécédent sur l'axe des abscisses.}
\end{exo}

\begin{exo}[Points et courbes]
\cas{1)}\par\nopagebreak
\explique{$f(a)=b$ signifie que le point de coordonnées $(a\,;b)$ appartient à $\Cf$.}{antécédent
en abscisse, image en ordonnée}
\explique{a) $f(0)=-3$ : le point $\res{(0\,;-3)}$ appartient à $\Cf$.}{}
\explique{b) $f(-4)=0$ : le point $\res{(-4\,;0)}$ appartient à $\Cf$.}{c'est un point de
l'axe des abscisses}
\explique{c) $f(2)=1$ et $f(6)=1$ : les points $\res{(2\,;1)}$ et $\res{(6\,;1)}$
appartiennent à $\Cf$.}{deux égalités, donc deux points, à la même hauteur}
\cas{2)}\par\nopagebreak
\explique{a) Le point $(3\,;-2)$ est sur $\Cf$ : $\res{f(3)=-2}$.}{l'abscisse est
l'antécédent, l'ordonnée l'image}
\explique{b) Un point de l'axe des ordonnées a pour abscisse $0$ : $\Cf$ passe par
$(0\,;4)$, donc $\res{f(0)=4}$.}{}
\explique{c) Un point de l'axe des abscisses a pour ordonnée $0$ : $\Cf$ passe par
$(-1\,;0)$ et $(5\,;0)$, donc $\res{f(-1)=0}$ et $\res{f(5)=0}$.}{}
\rem{Erreur fréquente en b) et c) : inverser les deux axes. L'axe des ordonnées est
l'axe vertical, où $x=0$.}
\cas{3) a)}\par\nopagebreak
\explique{Un point est sur la courbe lorsque l'image de son abscisse est égale à son
ordonnée : on calcule l'image de l'abscisse et on compare.}{méthode 12 du cours}
\begin{calculs}
f(-2) &= (-2)^{2}-3 &&\justif{l'image de l'abscisse de $A$}\\
f(-2) &= 4-3 = 1 &&
\end{calculs}
\explique{$f(-2)=1$, l'ordonnée de $A$ : $\res{A\in\Cf}$.}{}
\begin{calculs}
f(3) &= 3^{2}-3 &&\justif{l'image de l'abscisse de $B$}\\
f(3) &= 9-3 = 6 &&
\end{calculs}
\explique{$f(3)=6$, alors que l'ordonnée de $B$ est $5$ : $\res{B\notin\Cf}$.}{$6\ne 5$}
\begin{calculs}
f(1{,}5) &= 1{,}5^{2}-3 &&\justif{l'image de l'abscisse de $C$}\\
f(1{,}5) &= 2{,}25-3 = -0{,}75 &&\justif{$1{,}5\times 1{,}5=2{,}25$}
\end{calculs}
\explique{$f(1{,}5)=-0{,}75$, l'ordonnée de $C$ : $\res{C\in\Cf}$.}{}
\explique{Deux autres points : on choisit une abscisse et on calcule son image. $f(0)=-3$
et $f(1)=-2$ : les points $\res{(0\,;-3)}$ et $\res{(1\,;-2)}$ sont sur la
courbe.}{$0^{2}-3=-3$ ; $1^{2}-3=-2$}
\cas{3) b)}\par\nopagebreak
\begin{calculs}
f(-2) &= 3\times(-2)+1 &&\justif{l'image de l'abscisse de $A$}\\
f(-2) &= -6+1 = -5 &&
\end{calculs}
\explique{$f(-2)=-5$, l'ordonnée de $A$ : $\res{A\in\Cf}$.}{}
\begin{calculs}
f(2) &= 3\times 2+1 &&\justif{l'image de l'abscisse de $B$}\\
f(2) &= 7 &&
\end{calculs}
\explique{$f(2)=7$, alors que l'ordonnée de $B$ est $6$ : $\res{B\notin\Cf}$.}{$7\ne 6$}
\begin{calculs}
f\!\left(\tfrac13\right) &= 3\times\dfrac13+1 &&\justif{l'image de l'abscisse de $C$}\\
f\!\left(\tfrac13\right) &= 1+1 = 2 &&\justif{$3\times\frac13=\frac33=1$}
\end{calculs}
\explique{$f\!\left(\frac13\right)=2$, l'ordonnée de $C$ : $\res{C\in\Cf}$.}{}
\explique{Deux autres points : $f(0)=1$ et $f(1)=4$, d'où $\res{(0\,;1)}$ et
$\res{(1\,;4)}$.}{$3\times 0+1=1$ ; $3\times 1+1=4$}
\end{exo}

\begin{exo}[Trois antécédents]
\cas{a)}\par\nopagebreak
\begin{lecture}
Les nombres qui ont une image sont les abscisses des points de la courbe : elle va de
$-3$ à $3$. L'ensemble cherché est $\res{\intff{-3}{3}}$.

\smallskip
\gris{C'est l'ensemble de définition de $f$ ; extrémités marquées par des croix,
crochets fermés.}
\end{lecture}
\begin{figlecture}
\begin{tikzpicture}[x=0.4cm,y=0.4cm]
  \courbeXXVII
  \draw[black!70,dashed] (-3,-3)--(-3,0) (3,4)--(3,0);
  \morceau{-3}{3}\crochet{-3}{[}\crochet{3}{]}
\end{tikzpicture}
\end{figlecture}
\cas{b) et c)}\par\nopagebreak
\begin{lecture}
b) Au-dessus de $2$, la courbe est sur l'axe des abscisses : $f(2)=\res{0}$.

\medskip
c) La droite $y=-3$ ne rencontre la courbe qu'en son point de départ, d'abscisse $-3$ :
l'antécédent de $-3$ est $\res{-3}$.

\smallskip
\gris{Ensuite, la courbe remonte ; son point le plus bas, à droite, est à la hauteur
$-1$ : elle ne redescend jamais jusqu'à $-3$.}
\end{lecture}
\begin{figlecture}
\begin{tikzpicture}[x=0.4cm,y=0.4cm]
  \courbeXXVII
  \solx[coulCourbeC]{2}
  \droitey[coulInt]{-4}{4}{-3}
  \draw[coulInt,dashed] (-3,-3)--(-3,0);
  \solx{-3}
\end{tikzpicture}
\end{figlecture}
\cas{d)}\par\nopagebreak
\begin{lecture}
L'axe des abscisses, c'est-à-dire la droite $y=0$, coupe la courbe en trois points,
d'abscisses $-2\,;0$ et $2$ : le nombre $\res{0}$ a trois antécédents.

\smallskip
\gris{Tout nombre strictement compris entre $-1$ et $2$ convient aussi : la droite
horizontale coupe alors la courbe trois fois (en pointillés, la droite $y=1$).}
\end{lecture}
\begin{figlecture}
\begin{tikzpicture}[x=0.4cm,y=0.4cm]
  \courbeXXVII
  \solx{-2}\solx{0}\solx{2}
  \draw[black!60,dashed] (-4,1)--(4,1) node[above left,font=\scriptsize,inner sep=1pt]{$y=1$};
\end{tikzpicture}
\end{figlecture}
\cas{e)}\par\nopagebreak
\begin{lecture}
Le point le plus haut de la courbe a pour ordonnée $4$, le plus bas $-3$. La droite $y=5$
ne rencontre pas la courbe : $\res{5}$ n'a aucun antécédent.

\smallskip
\gris{Tout nombre strictement supérieur à $4$ ou strictement inférieur à $-3$ convient.}
\end{lecture}
\begin{figlecture}
\begin{tikzpicture}[x=0.4cm,y=0.4cm]
  \courbeXXVII
  \droitey[coulInt]{-4}{4}{5}
\end{tikzpicture}
\end{figlecture}
\end{exo}

\begin{exo}[Quelle courbe passe par ce point ?]
\explique{Un point est sur la courbe d'une fonction lorsque l'image de son abscisse est
égale à son ordonnée. On calcule donc l'image de l'abscisse par chaque
fonction.}{méthode 12 du cours}
\cas{1)}\par\nopagebreak
\explique{Pour $A(4\,;6)$, on calcule les images de $4$ et on les compare à $6$.}{}
\begin{calculs}
f(4) &= 10-4 = 6 &&\\
g(4) &= \dfrac{24}{4} = 6 &&\\
h(4) &= -2\times 4+14 = 6 &&\justif{$-8+14$}\\
k(4) &= -\dfrac12\times 4^{2}+5\times 4-4 &&\justif{la puissance d'abord : $4^{2}=16$}\\
k(4) &= -8+20-4 = 8 &&\justif{$-\frac12\times 16=-8$}
\end{calculs}
\explique{Les courbes de $\res{f}$, $\res{g}$ et $\res{h}$ passent par $A$ ; pas celle de
$k$, car $k(4)=8\ne 6$.}{un même point peut être sur plusieurs courbes}
\cas{2)}\par\nopagebreak
\explique{Pour $B(6\,;4)$ : images de $6$, comparées à $4$.}{}
\begin{calculs}
f(6) &= 10-6 = 4 &&\\
g(6) &= \dfrac{24}{6} = 4 &&\\
h(6) &= -2\times 6+14 = 2 &&\\
k(6) &= -\dfrac12\times 6^{2}+5\times 6-4 &&\justif{$6^{2}=36$}\\
k(6) &= -18+30-4 = 8 &&
\end{calculs}
\explique{Les courbes de $\res{f}$ et $\res{g}$ passent par $B$.}{$h(6)=2$ et $k(6)=8$ ne
valent pas $4$}
\explique{Pour $C(8\,;4)$ : images de $8$, comparées à $4$.}{}
\begin{calculs}
f(8) &= 10-8 = 2 &&\\
g(8) &= \dfrac{24}{8} = 3 &&\\
h(8) &= -2\times 8+14 = -2 &&\\
k(8) &= -\dfrac12\times 8^{2}+5\times 8-4 &&\justif{$8^{2}=64$}\\
k(8) &= -32+40-4 = 4 &&
\end{calculs}
\explique{Seule la courbe de $\res{k}$ passe par $C$.}{}

\smallskip
En résumé :
\begin{center}
\renewcommand{\arraystretch}{1.3}
\begin{tabular}{|c|c|c|c|c|}
\hline
 & $f$ & $g$ & $h$ & $k$\\\hline
image de $4$ & $\mathbf{6}$ & $\mathbf{6}$ & $\mathbf{6}$ & $8$\\\hline
image de $6$ & $\mathbf{4}$ & $\mathbf{4}$ & $2$ & $8$\\\hline
image de $8$ & $2$ & $3$ & $-2$ & $\mathbf{4}$\\\hline
\end{tabular}
\end{center}
\rem{Erreur fréquente : comparer l'image à l'abscisse du point. On calcule l'image de
l'\emph{abscisse}, et on la compare à l'\emph{ordonnée}.}
\end{exo}

\partie{Calculatrice}

\begin{exo}[Antécédents à la calculatrice]
\cas{1)}\par\nopagebreak
\explique{Application \texttt{Grapheur}, onglet \texttt{Expressions} : on entre
$f(x)=0{,}5x^{2}-2x+1$. Onglet \texttt{Tableau}, \texttt{Régler l'intervalle} :
\texttt{X début} $-2$, \texttt{X fin} $6$, \texttt{Pas} $1$.}{méthode 13 du cours :
début, fin et pas}
\begin{center}
\renewcommand{\arraystretch}{1.25}
\begin{tabular}{|c|*{9}{c|}}
\hline
$x$ & $-2$ & $-1$ & $0$ & $1$ & $2$ & $3$ & $4$ & $5$ & $6$\\\hline
$f(x)$ & $7$ & $3{,}5$ & $1$ & $-0{,}5$ & $-1$ & $-0{,}5$ & $1$ & $3{,}5$ & $7$\\\hline
\end{tabular}
\end{center}
\explique{Une valeur vérifiée à la main : $f(-2)=0{,}5\times 4+4+1$, soit $7$.}{$(-2)^{2}=4$
et $-2\times(-2)=4$}
\explique{Onglet \texttt{Graphique}, menu \texttt{Axes} : $X_{\min}=-2$, $X_{\max}=6$,
$Y_{\min}=-2$, $Y_{\max}=8$.}{un peu au-delà de la plus petite image ($-1$) et de la plus
grande ($7$)}

\cas{2) et 3)}\par\nopagebreak
\begin{lecture}
2) Sur la courbe : \texttt{OK}, \texttt{Calculer}, \texttt{Antécédent}, valeur $2$. La
calculatrice donne deux antécédents : $\res{x\approx -0{,}4}$ et $\res{x\approx 4{,}4}$.

\smallskip
\gris{Contrôle avec le tableau : $f(-1)=3{,}5$ et $f(0)=1$ ; $2$ est entre les deux, donc un
antécédent est entre $-1$ et $0$. De même, l'autre est entre $4$ et $5$.}

\medskip
3) Dans le tableau, $f(0)=1$ est positif et $f(1)=-0{,}5$ est négatif : la courbe traverse
l'axe des abscisses entre $0$ et $1$. Elle le retraverse entre $3$ et $4$. Le nombre $0$ a
donc \res{deux antécédents} : $x\approx 0{,}6$ et $x\approx 3{,}4$.

\smallskip
\gris{La plus petite image est $-1$, atteinte en $2$ : la courbe descend sous l'axe, puis
remonte.}
\end{lecture}
\begin{figlecture}
\begin{tikzpicture}[x=0.42cm,y=0.42cm]
  \repere{-3}{-2}{7}{8}
  \gradx{-2,-1,1,2,3,4,5,6}
  \grady{-1,1,2,3,4,5,6,7}
  \draw[coulCourbe,line width=1.1pt,domain=-2:6,samples=60] plot (\x,{0.5*(\x)*(\x)-2*(\x)+1});
  \foreach \a/\b in {-2/7,-1/3.5,0/1,1/-0.5,2/-1,3/-0.5,4/1,5/3.5,6/7} \path[coulCourbe] (\a,\b) node[croix]{};
  \node[coulCourbe,font=\footnotesize,left] at (5.8,7.4) {$\Cf$};
  \droitey[coulInt]{-3}{7}{2}
  \draw[coulInt,dashed] (-0.449,2)--(-0.449,0) (4.449,2)--(4.449,0);
  \solx{-0.449}\solx{4.449}
  \solx[coulIntB]{0.586}\solx[coulIntB]{3.414}
\end{tikzpicture}
\end{figlecture}
\rem{Les valeurs données par la calculatrice sont approchées : on écrit $\approx$, et on
arrondit comme demandé (à $0{,}1$).}
\end{exo}

\begin{exo}[Tableau avec un pas de 0,5]
\cas{1)}\par\nopagebreak
\explique{\texttt{Régler l'intervalle} : \texttt{X début} $-1$, \texttt{X fin} $2$,
\texttt{Pas} $0{,}5$.}{de $-1$ à $2$ en avançant de $0{,}5$ : sept valeurs}
\begin{center}
\renewcommand{\arraystretch}{1.25}
\begin{tabular}{|c|*{7}{c|}}
\hline
$x$ & $-1$ & $-0{,}5$ & $0$ & $0{,}5$ & $1$ & $1{,}5$ & $2$\\\hline
$g(x)$ & $0$ & $2$ & $3$ & $3$ & $2$ & $0$ & $-3$\\\hline
\end{tabular}
\end{center}
\explique{Une valeur vérifiée à la main :}{pour contrôler la saisie de la formule}
\begin{calculs}
g(-0{,}5) &= -2\times(-0{,}5)^{2}+(-0{,}5)+3 &&\justif{$-0{,}5$ entre parenthèses}\\
g(-0{,}5) &= -2\times 0{,}25-0{,}5+3 &&\justif{$(-0{,}5)^{2}=0{,}25$}\\
g(-0{,}5) &= -0{,}5-0{,}5+3 = 2 &&
\end{calculs}

\cas{2) et 3)}\par\nopagebreak
\begin{lecture}
2) \texttt{Axes} : $X_{\min}=-1$, $X_{\max}=2$, $Y_{\min}=-4$, $Y_{\max}=4$. Sur la copie,
on place les sept points du tableau et on les relie par une courbe sans angle.

\medskip
3) \texttt{Aller à} $0{,}25$ : $g(0{,}25)=\res{3{,}125}$.

\smallskip
\gris{$0{,}25$ n'est pas dans le tableau ; à la main :
$-2\times 0{,}0625+0{,}25+3=3{,}125$.}

\medskip
Antécédents de $2$ (en rouge) : $\res{-0{,}5}$ et $\res{1}$, lus dans le tableau ; la
droite $y=2$ coupe la courbe en ces deux points.

Antécédents de $0$ (en bleu) : $\res{-1}$ et $\res{1{,}5}$.

\smallskip
\gris{Ces valeurs sont exactes : elles figurent dans le tableau.}
\end{lecture}
\begin{figlecture}
\begin{tikzpicture}[x=0.9cm,y=0.9cm]
  \repere{-2}{-4}{3}{4}
  \gradx{-1,1,2}
  \grady{-3,-2,-1,1,2,3}
  \draw[coulCourbe,line width=1.1pt,domain=-1:2,samples=60] plot (\x,{-2*(\x)*(\x)+(\x)+3});
  \foreach \a/\b in {-1/0,-0.5/2,0/3,0.5/3,1/2,1.5/0,2/-3} \path[coulCourbe] (\a,\b) node[croix]{};
  \node[coulCourbe,font=\footnotesize,right] at (1.6,-1.6) {$\mathcal{C}_{g}$};
  \droitey[coulInt]{-2}{3}{2}
  \draw[coulInt,dashed] (-0.5,2)--(-0.5,0) (1,2)--(1,0);
  \solx{-0.5}\solx{1}
  \solx[coulIntB]{-1}\solx[coulIntB]{1.5}
\end{tikzpicture}
\end{figlecture}
\end{exo}

\begin{exo}[Une fonction quotient]
\cas{1)}\par\nopagebreak
\explique{On entre $f(x)=\dfrac{2x-1}{(x+1)^{2}}$, avec les parenthèses autour de $x+1$.
\texttt{Régler l'intervalle} : \texttt{X début} $0$, \texttt{X fin} $5$, \texttt{Pas}
$0{,}5$.}{sans parenthèses, la calculatrice élève au carré le $1$ seulement}
\begin{center}
\renewcommand{\arraystretch}{1.25}
\begin{tabular}{|c|*{6}{c|}}
\hline
$x$ & $0$ & $0{,}5$ & $1$ & $1{,}5$ & $2$ & $2{,}5$\\\hline
$f(x)$ & $-1$ & $0$ & $0{,}25$ & $0{,}32$ & $0{,}33$ & $0{,}33$\\\hline
\end{tabular}

\smallskip
\begin{tabular}{|c|*{5}{c|}}
\hline
$x$ & $3$ & $3{,}5$ & $4$ & $4{,}5$ & $5$\\\hline
$f(x)$ & $0{,}31$ & $0{,}30$ & $0{,}28$ & $0{,}26$ & $0{,}25$\\\hline
\end{tabular}
\end{center}
\explique{Une valeur vérifiée à la main : $f(1)=\dfrac{2-1}{2^{2}}$, soit
$\dfrac14=0{,}25$.}{valeurs arrondies à $0{,}01$ : $f(2)=\frac39=0{,}333\ldots$}

\cas{2)}\par\nopagebreak
\explique{\texttt{Axes} : $X_{\min}=0$, $X_{\max}=5$, $Y_{\min}=-1{,}2$,
$Y_{\max}=0{,}5$.}{la plus petite image est $-1$, la plus grande environ $0{,}33$}

\cas{3) et 4)}\par\nopagebreak
\begin{center}
\begin{tikzpicture}[x=2.2cm,y=2.2cm]
  \draw[-{Stealth[length=2mm]},thick] (-0.2,0)--(5.4,0) node[below left=-1pt]{$x$};
  \draw[-{Stealth[length=2mm]},thick] (0,-1.15)--(0,0.65) node[below left=-1pt]{$y$};
  \foreach \k in {0.5,1,...,5} \draw (\k,0.03)--(\k,-0.03);
  \foreach \k in {1,2,3,4,5} \node[below,font=\scriptsize,inner sep=2pt] at (\k,-0.03) {$\k$};
  \node[below,font=\scriptsize,inner sep=2pt] at (0.5,-0.03) {$0{,}5$};
  \foreach \k/\l in {-1/-1,-0.5/-0{,}5,0.25/0{,}25,0.5/0{,}5} {\draw (0.03,\k)--(-0.03,\k) node[left,font=\scriptsize,inner sep=2pt]{$\l$};}
  \node[below left,font=\scriptsize,inner sep=1.5pt] at (0,0) {$0$};
  \draw[coulCourbe,line width=1.1pt,domain=0:5,samples=100] plot (\x,{(2*(\x)-1)/(((\x)+1)*((\x)+1))});
  \node[coulCourbe,font=\footnotesize] at (2.5,0.45) {$\Cf$};
  \draw[black!70,thick] (0,0.25)--(5.3,0.25) node[above left,font=\scriptsize,inner sep=1pt]{$y=0{,}25$};
  \draw[black!70,dashed] (1,0.25)--(1,0) (5,0.25)--(5,0);
  \morceau{1}{5}\crochet{1}{[}\crochet{5}{]}
  \path[coulIntB] (0.5,0) node[croix]{};
\end{tikzpicture}
\end{center}
\explique{3) Le tableau donne $f(1)=0{,}25$ et $f(5)=0{,}25$ ; la droite $y=0{,}25$ coupe
la courbe en ces deux points seulement. Les antécédents de $0{,}25$ sont $\res{1}$ et
$\res{5}$.}{\texttt{Calculer}, \texttt{Antécédent} confirme ; ce sont des valeurs exactes}
\begin{calculs}
f(5) &= \dfrac{2\times 5-1}{(5+1)^{2}} &&\justif{on vérifie le second à la main}\\
f(5) &= \dfrac{9}{36} &&\\
f(5) &= \dfrac{1\times\barre{9}}{4\times\barre{9}} = \dfrac14 &&\justif{$0{,}25=\frac14$}
\end{calculs}
\explique{4) La courbe coupe l'axe des abscisses en $0{,}5$ (croix bleue) :
$\res{S=\ens{0{,}5}}$.}{un quotient est nul lorsque son numérateur est nul : $2x-1=0$
pour $x=0{,}5$}
\explique{La courbe est sur la droite $y=0{,}25$ ou au-dessus de $1$ à $5$ (en rouge sur
l'axe) : $\res{S=\intff{1}{5}}$.}{inégalité large : les bornes $1$ et $5$ sont
comprises}

\cas{5)}\par\nopagebreak
\begin{center}
\begin{tikzpicture}[x=2.2cm,y=2.2cm]
  \draw[-{Stealth[length=2mm]},thick] (-0.2,0)--(5.4,0) node[below left=-1pt]{$x$};
  \draw[-{Stealth[length=2mm]},thick] (0,-1.15)--(0,0.65) node[below left=-1pt]{$y$};
  \foreach \k in {0.5,1,...,5} \draw (\k,0.03)--(\k,-0.03);
  \foreach \k in {2,3,4,5} \node[below,font=\scriptsize,inner sep=2pt] at (\k,-0.03) {$\k$};
  \node[below,font=\scriptsize,inner sep=2pt] at (0.95,-0.03) {$1$};
  \foreach \k/\l in {-1/-1,-0.5/-0{,}5} {\draw (0.03,\k)--(-0.03,\k) node[left,font=\scriptsize,inner sep=2pt]{$\l$};}
  \node[below left,font=\scriptsize,inner sep=1.5pt] at (0,0) {$0$};
  \draw[coulCourbe,line width=1.1pt,domain=0:5,samples=100] plot (\x,{(2*(\x)-1)/(((\x)+1)*((\x)+1))});
  \node[coulCourbe,font=\footnotesize] at (2.5,0.45) {$\Cf$};
  \draw[coulCourbeC,line width=1.1pt] (0,-1)--(1.6,0.6) node[right,font=\footnotesize]{$y=x-1$};
  \path (0,-1) node[croix]{} (1.303,0.303) node[croix]{};
  \draw[black!70,dashed] (1.303,0.303)--(1.303,0);
  \solx{0}\solx{1.303}
  \node[coulInt,below,font=\scriptsize,inner sep=2pt] at (1.303,-0.06) {$\approx 1{,}30$};
\end{tikzpicture}
\end{center}
\explique{On entre aussi $g(x)=x-1$ ; sur une courbe, \texttt{Calculer},
\texttt{Intersection}. Les deux courbes se coupent en deux points, d'abscisses $0$ et
environ $1{,}30$ : $\res{S\approx\ens{0\,;1{,}30}}$.}{les solutions de l'équation sont les
abscisses des points communs}
\explique{$0$ est une solution exacte : $f(0)=\dfrac{-1}{1}=-1$ et $0-1=-1$.}{l'autre n'est
connue qu'arrondie}
\rem{La courbe de $f$ varie très peu entre $1$ et $5$ : à l'écran, on peut resserrer
l'axe des ordonnées pour bien voir qu'elle reste au-dessus de la droite $y=0{,}25$.}
\end{exo}

\begin{exo}[Intersections de courbes]
\cas{1)}\par\nopagebreak
\begin{lecture}
On entre $f(x)=x^{3}-2x^{2}$ et $g(x)=x-2$, puis \texttt{Axes} : $X_{\min}=-2$,
$X_{\max}=3$ ; pour bien voir les points communs, $Y_{\min}=-5$ et $Y_{\max}=4$.

\smallskip
\gris{$f(-2)=-16$ : avec toute la courbe, les points communs seraient écrasés.}

\medskip
\texttt{Calculer}, \texttt{Intersection} : les courbes se coupent aux points d'abscisses
$-1\,;1$ et $2$.

\smallskip
\gris{Vérification : $f(-1)=-1-2=-3$ et $g(-1)=-3$ ; $f(1)=1-2=-1$ et $g(1)=-1$ ;
$f(2)=8-8=0$ et $g(2)=0$.}

\medskip
$x^{3}-2x^{2}\le x-2$ : on cherche où $\Cf$ est au-dessous de $\Cg$, ou la coupe : de $-2$
à $-1$ (en rouge), et de $1$ à $2$ (en bleu).

$\res{S=\intff{-2}{-1}\cup\intff{1}{2}}$.

\smallskip
\gris{Inégalité large : les abscisses des points communs sont comprises.}
\end{lecture}
\begin{figlecture}
\begin{tikzpicture}[x=0.55cm,y=0.55cm]
  \repere{-3}{-5}{4}{4}
  \gradx{-2,-1,1,2,3}
  \grady{-4,-3,-2,-1,1,2,3}
  \begin{scope}
    \clip (-3,-5) rectangle (4,4);
    \draw[coulCourbe,line width=1.1pt,domain=-2:3,samples=80] plot (\x,{(\x)*(\x)*(\x)-2*(\x)*(\x)});
  \end{scope}
  \draw[coulCourbeC,line width=1.1pt] (-2,-4)--(3,1);
  \node[coulCourbe,font=\footnotesize,left] at (2.6,3.6) {$\Cf$};
  \node[coulCourbeC,font=\footnotesize,right] at (3,1) {$\Cg$};
  \path (-1,-3) node[croix]{} (1,-1) node[croix]{} (2,0) node[croix]{};
  \draw[black!70,dashed] (-1,-3)--(-1,0) (1,-1)--(1,0);
  \morceau{-2}{-1}\crochet{-2}{[}\crochet{-1}{]}
  \morceau[coulIntB]{1}{2}\crochet[coulIntB]{1}{[}\crochet[coulIntB]{2}{]}
\end{tikzpicture}
\end{figlecture}

\cas{2)}\par\nopagebreak
\explique{La calculatrice (\texttt{Intersection}) donne un seul point commun,
$\res{(1\,;9)}$. Par le calcul : un point commun a la même ordonnée sur les deux
courbes.}{les deux expressions de $y$ sont égales}
\begin{calculs}
3x^{2}+x+5 &= 3x^{2}-2x+8 &&\justif{les deux ordonnées sont égales}\\
x+5 &= -2x+8 &&\justif{on retranche $3x^{2}$ aux deux membres}\\
3x+5 &= 8 &&\justif{on ajoute $2x$}\\
3x &= 3 &&\justif{on retranche $5$}\\
x &= 1 &&\justif{on divise par $3$}
\end{calculs}
\begin{calculs}
y &= 3\times 1^{2}+1+5 &&\justif{on calcule l'ordonnée avec la première équation}\\
y &= 9 &&\justif{avec la seconde : $3-2+8=9$, même résultat}
\end{calculs}
\explique{Les courbes ont un seul point commun, $\res{(1\,;9)}$ : le calcul confirme la
calculatrice.}{les $3x^{2}$ disparaissent : il reste une équation du premier degré}

\cas{3)}\par\nopagebreak
\begin{lecture}
À l'écran, les deux courbes ne se rencontrent pas. Par le calcul, pour $x\ne 0$ :

\smallskip
$\dfrac1x=\dfrac{1+2x}{x}$ \ donne \ $1=1+2x$ \ (même dénominateur), \ puis \ $2x=0$,
\ soit \ $x=0$.

\smallskip
Or $0$ est une valeur interdite des deux fonctions : \res{aucun point commun}.

\smallskip
\gris{On peut le voir autrement : $\dfrac{1+2x}{x}=\dfrac1x+\dfrac{2x}{x}=\dfrac1x+2$. La
seconde courbe est la première montée de $2$ : elles ne se coupent jamais.}
\end{lecture}
\begin{figlecture}
\begin{tikzpicture}[x=0.5cm,y=0.5cm]
  \repere{-4}{-3}{4}{6}
  \gradx{-3,-2,-1,1,2,3}
  \grady{-2,-1,1,2,3,4,5}
  \begin{scope}
    \clip (-4,-3) rectangle (4,6);
    \draw[coulCourbe,line width=1.1pt,domain=0.15:4,samples=60] plot (\x,{1/(\x)});
    \draw[coulCourbe,line width=1.1pt,domain=-4:-0.15,samples=60] plot (\x,{1/(\x)});
    \draw[coulCourbeC,line width=1.1pt,domain=0.22:4,samples=60] plot (\x,{2+1/(\x)});
    \draw[coulCourbeC,line width=1.1pt,domain=-4:-0.15,samples=60] plot (\x,{2+1/(\x)});
  \end{scope}
  \node[coulCourbe,font=\footnotesize] at (-2.4,-1.5) {$y=\frac1x$};
  \node[coulCourbeC,font=\footnotesize] at (2.4,3.6) {$y=\frac{1+2x}{x}$};
\end{tikzpicture}
\end{figlecture}
\end{exo}

\partie{Résolution graphique}

\begin{exo}[Équations]
\explique{Les solutions de $f(x)=k$ sont les abscisses des points de $\Cf$ situés sur la
droite horizontale $y=k$. Une équation se conclut par des accolades.}{propriété 3 du
cours : ce sont les antécédents de $k$}
\cas{1) a)}\par\nopagebreak
\begin{lecture}
La droite $y=2$ touche $\Cf$ en un seul point, son point le plus haut, d'abscisse $5$ :
$\res{S=\ens{5}}$.
\end{lecture}
\begin{figlecture}
\begin{tikzpicture}[x=0.42cm,y=0.42cm]
  \courbeXXXIII
  \droitey[coulInt]{-1}{11}{2}
  \draw[coulInt,dashed] (5,2)--(5,0);
  \solx{5}
\end{tikzpicture}
\end{figlecture}
\cas{1) b)}\par\nopagebreak
\begin{lecture}
La droite $y=1$ coupe $\Cf$ en deux points, d'abscisses $4$ et $6$ :
$\res{S=\ens{4\,;6}}$.

\smallskip
\gris{Vérification : les points $(4\,;1)$ et $(6\,;1)$ sont sur la courbe.}
\end{lecture}
\begin{figlecture}
\begin{tikzpicture}[x=0.42cm,y=0.42cm]
  \courbeXXXIII
  \droitey[coulInt]{-1}{11}{1}
  \draw[coulInt,dashed] (4,1)--(4,0) (6,1)--(6,0);
  \solx{4}\solx{6}
\end{tikzpicture}
\end{figlecture}
\cas{1) c)}\par\nopagebreak
\begin{lecture}
La droite $y=0$ est l'axe des abscisses : $\Cf$ le coupe en $0\,;3$ et $8$.
$\res{S=\ens{0\,;3\,;8}}$.

\smallskip
\gris{La courbe part du point $(0\,;0)$ : on n'oublie pas la solution $0$.}
\end{lecture}
\begin{figlecture}
\begin{tikzpicture}[x=0.42cm,y=0.42cm]
  \courbeXXXIII
  \solx{0}\solx{3}\solx{8}
\end{tikzpicture}
\end{figlecture}
\cas{1) d)}\par\nopagebreak
\begin{lecture}
La droite $y=-2$ touche $\Cf$ en son point le plus bas, d'abscisse $1$ :
$\res{S=\ens{1}}$.
\end{lecture}
\begin{figlecture}
\begin{tikzpicture}[x=0.42cm,y=0.42cm]
  \courbeXXXIII
  \droitey[coulInt]{-1}{11}{-2}
  \draw[coulInt,dashed] (1,-2)--(1,0);
  \solx{1}
\end{tikzpicture}
\end{figlecture}
\cas{1) e)}\par\nopagebreak
\begin{lecture}
La droite $y=3$ passe au-dessus de toute la courbe, qui ne dépasse pas $2$ :
$\res{S=\varnothing}$.

\smallskip
\gris{Pas de point commun : l'équation n'a pas de solution.}
\end{lecture}
\begin{figlecture}
\begin{tikzpicture}[x=0.42cm,y=0.42cm]
  \courbeXXXIII
  \droitey[coulInt]{-1}{11}{3}
\end{tikzpicture}
\end{figlecture}
\cas{2)}\par\nopagebreak
\begin{lecture}
La droite $y=-1$ coupe $\Cf$ en \res{trois points} : un entre $0$ et $1$, un en $2$, un en
$10$. L'équation $f(x)=-1$ a trois solutions.

\smallskip
\gris{On ne demande que le nombre de solutions : inutile de lire la première, qui n'est
pas un nombre entier.}
\end{lecture}
\begin{figlecture}
\begin{tikzpicture}[x=0.42cm,y=0.42cm]
  \courbeXXXIII
  \droitey[coulInt]{-1}{11}{-1}
  \draw[coulInt,dashed] (2,-1)--(2,0) (10,-1)--(10,0);
  \path[coulInt] (0.4,-1) node[croix]{} (2,-1) node[croix]{} (10,-1) node[croix]{};
\end{tikzpicture}
\end{figlecture}
\end{exo}

\begin{exo}[Inéquations]
\explique{Les solutions de $f(x)>k$ sont les abscisses des points de $\Cf$ strictement
au-dessus de la droite $y=k$ ; pour $f(x)<k$, strictement au-dessous. Pour $\ge$ ou
$\le$, on garde aussi les points sur la droite. Une inéquation se conclut par des
intervalles.}{propriété 3 du cours}
\cas{a)}\par\nopagebreak
\begin{lecture}
$\Cf$ est strictement au-dessus de la droite $y=-2$ partout, sauf en $1$, où elle la
touche :

$\res{S=\intfo{0}{1}\cup\intof{1}{10}}$.

\smallskip
\gris{$f(1)=-2$ n'est pas strictement supérieur à $-2$ : on retire $1$.}
\end{lecture}
\begin{figlecture}
\begin{tikzpicture}[x=0.42cm,y=0.42cm]
  \courbeXXXIII
  \droitey[black!70]{-1}{11}{-2}
  \draw[black!70,dashed] (1,-2)--(1,0);
  \morceau{0}{1}\crochet{0}{[}\crochet{1}{[}
  \morceau[coulIntB]{1}{10}\crochet[coulIntB]{1}{]}\crochet[coulIntB]{10}{]}
\end{tikzpicture}
\end{figlecture}
\cas{b)}\par\nopagebreak
\begin{lecture}
$\Cf$ est strictement au-dessous de la droite $y=1$ de $0$ à $4$, puis de $6$ à $10$ :

$\res{S=\intfo{0}{4}\cup\intof{6}{10}}$.

\smallskip
\gris{En $4$ et en $6$, $f(x)=1$ : exclus, car l'inégalité est stricte. En $0$ et $10$,
$f(x)<1$ : inclus.}
\end{lecture}
\begin{figlecture}
\begin{tikzpicture}[x=0.42cm,y=0.42cm]
  \courbeXXXIII
  \droitey[black!70]{-1}{11}{1}
  \draw[black!70,dashed] (4,1)--(4,0) (6,1)--(6,0);
  \morceau{0}{4}\crochet{0}{[}\crochet{4}{[}
  \morceau[coulIntB]{6}{10}\crochet[coulIntB]{6}{]}\crochet[coulIntB]{10}{]}
\end{tikzpicture}
\end{figlecture}
\cas{c)}\par\nopagebreak
\begin{lecture}
La courbe ne monte jamais au-dessus de la droite $y=2$ : elle est au-dessous de la droite ou sur elle
sur tout l'intervalle de définition.

$\res{S=\intff{0}{10}}$.
\end{lecture}
\begin{figlecture}
\begin{tikzpicture}[x=0.42cm,y=0.42cm]
  \courbeXXXIII
  \droitey[black!70]{-1}{11}{2}
  \morceau{0}{10}\crochet{0}{[}\crochet{10}{]}
\end{tikzpicture}
\end{figlecture}
\cas{d)}\par\nopagebreak
\begin{lecture}
$\Cf$ est sur la droite $y=1$ ou au-dessus de $4$ à $6$ :

$\res{S=\intff{4}{6}}$.

\smallskip
\gris{Inégalité large : $4$ et $6$ sont compris.}
\end{lecture}
\begin{figlecture}
\begin{tikzpicture}[x=0.42cm,y=0.42cm]
  \courbeXXXIII
  \droitey[black!70]{-1}{11}{1}
  \draw[black!70,dashed] (4,1)--(4,0) (6,1)--(6,0);
  \morceau{4}{6}\crochet{4}{[}\crochet{6}{]}
\end{tikzpicture}
\end{figlecture}
\rem{On ne conclut que sur la partie tracée de la courbe : les solutions restent dans
$\intff{0}{10}$, l'ensemble de définition.}
\end{exo}

\begin{exo}[Équation et inéquations]
\cas{a)}\par\nopagebreak
\begin{lecture}
La droite $y=1$ coupe $\mathcal{C}$ en trois points, d'abscisses $-4\,;0$ et $4$ :

$\res{S=\ens{-4\,;0\,;4}}$.

\smallskip
\gris{Le point de départ $(-4\,;1)$ compte : $-4$ est dans l'ensemble de définition.}
\end{lecture}
\begin{figlecture}
\begin{tikzpicture}[x=0.45cm,y=0.45cm]
  \courbeXXXV
  \droitey[coulInt]{-5}{6}{1}
  \draw[coulInt,dashed] (-4,1)--(-4,0) (4,1)--(4,0);
  \solx{-4}\solx{0}\solx{4}
\end{tikzpicture}
\end{figlecture}
\cas{b)}\par\nopagebreak
\begin{lecture}
$\mathcal{C}$ est sur la droite $y=1$ ou au-dessous de $-4$ à $0$, puis de $4$ à $5$ :

$\res{S=\intff{-4}{0}\cup\intff{4}{5}}$.

\smallskip
\gris{Inégalité large : les abscisses des points sur la droite sont comprises.}
\end{lecture}
\begin{figlecture}
\begin{tikzpicture}[x=0.45cm,y=0.45cm]
  \courbeXXXV
  \droitey[black!70]{-5}{6}{1}
  \draw[black!70,dashed] (-4,1)--(-4,0) (4,1)--(4,0);
  \morceau{-4}{0}\crochet{-4}{[}\crochet{0}{]}
  \morceau[coulIntB]{4}{5}\crochet[coulIntB]{4}{[}\crochet[coulIntB]{5}{]}
\end{tikzpicture}
\end{figlecture}
\cas{c)}\par\nopagebreak
\begin{lecture}
$\mathcal{C}$ est strictement au-dessous de la droite $y=-1$ entre $-3$ et $-1$ :

$\res{S=\intoo{-3}{-1}}$.

\smallskip
\gris{En $-3$ et en $-1$, $f(x)=-1$ : exclus, car l'inégalité est stricte.}
\end{lecture}
\begin{figlecture}
\begin{tikzpicture}[x=0.45cm,y=0.45cm]
  \courbeXXXV
  \droitey[black!70]{-5}{6}{-1}
  \draw[black!70,dashed] (-3,-1)--(-3,0) (-1,-1)--(-1,0);
  \morceau{-3}{-1}\crochet{-3}{]}\crochet{-1}{[}
\end{tikzpicture}
\end{figlecture}
\end{exo}

\begin{exo}[Avec la formule]
\cas{1)}\par\nopagebreak
\explique{On entre $f(x)=\dfrac{x^{3}-7x+6}{2}$ ; \texttt{Régler l'intervalle} :
\texttt{X début} $-3$, \texttt{X fin} $3$, \texttt{Pas} $1$.}{numérateur entre
parenthèses si on le tape sur une ligne : $(x^{3}-7x+6)/2$}
\begin{center}
\renewcommand{\arraystretch}{1.25}
\begin{tabular}{|c|*{7}{c|}}
\hline
$x$ & $-3$ & $-2$ & $-1$ & $0$ & $1$ & $2$ & $3$\\\hline
$f(x)$ & $0$ & $6$ & $6$ & $3$ & $0$ & $0$ & $6$\\\hline
\end{tabular}
\end{center}
\explique{Une valeur vérifiée à la main :}{}
\begin{calculs}
f(-2) &= \dfrac{(-2)^{3}-7\times(-2)+6}{2} &&\justif{$-2$ entre parenthèses}\\
f(-2) &= \dfrac{-8+14+6}{2} &&\justif{$(-2)^{3}=-8$ et $-7\times(-2)=14$}\\
f(-2) &= \dfrac{12}{2} = 6 &&
\end{calculs}

\cas{2)}\par\nopagebreak
\begin{lecture}
La courbe coupe l'axe des ordonnées au point d'ordonnée $3$ : $f(0)=\res{3}$.

\medskip
La droite $y=6$ coupe la courbe aux points d'abscisses $-2\,;-1$ et $3$ :

$\res{S=\ens{-2\,;-1\,;3}}$.

\smallskip
\gris{Entre $-2$ et $-1$, la courbe passe un peu au-dessus de $6$, puis redescend.}
\end{lecture}
\begin{figlecture}
\begin{tikzpicture}[x=0.5cm,y=0.5cm]
  \courbeXXXVI
  \path[coulIntB] (0,3) node[croix]{};
  \droitey[coulInt]{-4}{4}{6}
  \draw[coulInt,dashed] (-2,6)--(-2,0) (-1,6)--(-1,0) (3,6)--(3,0);
  \solx{-2}\solx{-1}\solx{3}
\end{tikzpicture}
\end{figlecture}

\cas{3)}\par\nopagebreak
\explique{On développe le produit, deux facteurs à la fois.}{on part du membre de droite,
le plus compliqué}
\begin{calculs}
(x+1)(x+2)(x-3) &= (x^{2}+2x+x+2)(x-3) &&\justif{on développe les deux premiers facteurs}\\
(x+1)(x+2)(x-3) &= (x^{2}+3x+2)(x-3) &&\justif{$2x+x=3x$}\\
(x+1)(x+2)(x-3) &= x^{3}-3x^{2}+3x^{2}-9x+2x-6 &&\justif{chaque terme multiplie $x$, puis $-3$}\\
(x+1)(x+2)(x-3) &= x^{3}-7x-6 &&\justif{$-3x^{2}+3x^{2}=0$ ; $-9x+2x=-7x$}
\end{calculs}
\explique{On résout ensuite $f(x)=6$ par le calcul :}{}
\begin{calculs}
\dfrac{x^{3}-7x+6}{2} &= 6 &&\justif{on recopie l'équation}\\
x^{3}-7x+6 &= 12 &&\justif{on multiplie les deux membres par $2$}\\
x^{3}-7x-6 &= 0 &&\justif{on retranche $12$}\\
(x+1)(x+2)(x-3) &= 0 &&\justif{la forme factorisée vérifiée au-dessus}
\end{calculs}
\explique{Un produit est nul lorsque l'un de ses facteurs est nul : $x+1=0$ ou $x+2=0$ ou
$x-3=0$. Les solutions exactes sont $\res{-2\,;-1}$ et $\res{3}$.}{rappel du collège ; le calcul
confirme la lecture graphique de la question 2}

\cas{4)}\par\nopagebreak
\begin{lecture}
$f(x)\ge 6$ : la courbe est sur la droite $y=6$ ou au-dessus de $-2$ à $-1$ (en rouge), et
en $3$ seulement, son extrémité droite (croix bleue) :

$\res{S=\intff{-2}{-1}\cup\ens{3}}$.

\smallskip
\gris{En $3$, un seul nombre : on l'écrit entre accolades.}

\medskip
$f(x)<0$ : la courbe est strictement au-dessous de l'axe des abscisses entre $1$ et $2$
(en vert) :

$\res{S=\intoo{1}{2}}$.

\smallskip
\gris{$f(1)=0$ et $f(2)=0$ ne sont pas strictement négatifs ; en $-3$, $f(-3)=0$ non plus.}
\end{lecture}
\begin{figlecture}
\begin{tikzpicture}[x=0.5cm,y=0.5cm]
  \courbeXXXVI
  \droitey[black!70]{-4}{4}{6}
  \draw[black!70,dashed] (-2,6)--(-2,0) (-1,6)--(-1,0) (3,6)--(3,0);
  \morceau{-2}{-1}\crochet{-2}{[}\crochet{-1}{]}
  \solx[coulIntB]{3}
  \morceau[coulIntC]{1}{2}\crochet[coulIntC]{1}{]}\crochet[coulIntC]{2}{[}
\end{tikzpicture}
\end{figlecture}

\cas{5)}\par\nopagebreak
\begin{lecture}
$f(x)=0$ : la courbe coupe l'axe des abscisses en $-3\,;1$ et $2$ : \res{trois
solutions}.

\medskip
$f(x)=4$ : la droite $y=4$ coupe la courbe trois fois (en rouge) : \res{trois solutions}.

\smallskip
\gris{Une entre $-3$ et $-2$ ($f$ passe de $0$ à $6$), une entre $-1$ et $0$ (de $6$ à
$3$), une entre $2$ et $3$ (de $0$ à $6$).}

\medskip
$f(x)=-1$ : la courbe descend à peine sous l'axe entre $1$ et $2$, sans atteindre $-1$ ;
la droite $y=-1$ (en bleu) ne la rencontre pas : \res{aucune solution}.
\end{lecture}
\begin{figlecture}
\begin{tikzpicture}[x=0.5cm,y=0.5cm]
  \courbeXXXVI
  \solx[coulCourbeC]{-3}\solx[coulCourbeC]{1}\solx[coulCourbeC]{2}
  \droitey[coulInt]{-4}{4}{4}
  \path[coulInt] (-2.49,4) node[croix]{} (-0.29,4) node[croix]{} (2.78,4) node[croix]{};
  \droitey[coulIntB]{-4}{4}{-1}
\end{tikzpicture}
\end{figlecture}
\rem{Pour un nombre de solutions, une lecture approchée suffit : on compte les points
communs entre la courbe et la droite horizontale.}
\end{exo}

\partie{Comparer deux fonctions}

\begin{exo}[Deux courbes]
\explique{Les solutions de $f(x)=g(x)$ sont les abscisses des points d'intersection de
$\Cf$ et $\Cg$ ; celles de $f(x)>g(x)$, les abscisses des points où $\Cf$ est strictement
au-dessus de $\Cg$.}{propriété 4 du cours}
\cas{a)}\par\nopagebreak
\begin{lecture}
Les deux courbes se coupent aux points $(-2\,;2)$, $(2\,;0)$ et $(6\,;2)$ :

$\res{S=\ens{-2\,;2\,;6}}$.

\smallskip
\gris{On lit les abscisses des points communs, pas leurs ordonnées.}
\end{lecture}
\begin{figlecture}
\begin{tikzpicture}[x=0.42cm,y=0.42cm]
  \courbesXXXVII
  \path (-2,2) node[croix]{} (2,0) node[croix]{} (6,2) node[croix]{};
  \draw[black!70,dashed] (-2,2)--(-2,0) (6,2)--(6,0);
  \solx{-2}\solx{2}\solx{6}
\end{tikzpicture}
\end{figlecture}
\cas{b)}\par\nopagebreak
\begin{lecture}
On regarde quelle courbe est au-dessus entre deux points d'intersection qui se suivent.
$\Cf$ est strictement au-dessus de $\Cg$ entre $-2$ et $2$, puis après $6$ :

$\res{S=\intoo{-2}{2}\cup\intof{6}{7}}$.

\smallskip
\gris{Inégalité stricte : les abscisses des points communs sont exclues ; $7$ est
inclus, car $f(7)=3$ et $g(7)=1$.}
\end{lecture}
\begin{figlecture}
\begin{tikzpicture}[x=0.42cm,y=0.42cm]
  \courbesXXXVII
  \draw[black!70,dashed] (-2,2)--(-2,0) (6,2)--(6,0);
  \morceau{-2}{2}\crochet{-2}{]}\crochet{2}{[}
  \morceau[coulIntB]{6}{7}\crochet[coulIntB]{6}{]}\crochet[coulIntB]{7}{]}
\end{tikzpicture}
\end{figlecture}
\cas{c)}\par\nopagebreak
\begin{lecture}
$\Cf$ est au-dessous de $\Cg$, ou la coupe, de $-3$ à $-2$ et de $2$ à $6$ :

$\res{S=\intff{-3}{-2}\cup\intff{2}{6}}$.

\smallskip
\gris{C'est le reste de l'intervalle de définition $\intff{-3}{7}$ ; inégalité large : les
abscisses des points communs sont comprises.}
\end{lecture}
\begin{figlecture}
\begin{tikzpicture}[x=0.42cm,y=0.42cm]
  \courbesXXXVII
  \draw[black!70,dashed] (-2,2)--(-2,0) (6,2)--(6,0);
  \morceau{-3}{-2}\crochet{-3}{[}\crochet{-2}{]}
  \morceau[coulIntB]{2}{6}\crochet[coulIntB]{2}{[}\crochet[coulIntB]{6}{]}
\end{tikzpicture}
\end{figlecture}
\rem{Erreur fréquente : comparer les courbes à l'axe des abscisses. Ici, on compare
$\Cf$ à $\Cg$, pas à la droite $y=0$.}
\end{exo}

\begin{exo}[Une droite et une courbe]
\cas{a)}\par\nopagebreak
\begin{lecture}
$f(x)\le 2$ : on compare la droite $\Cf$ à la droite horizontale $y=2$. Elle la coupe en
$0$, et elle est au-dessous à droite de $0$ :

$\res{S=\intff{0}{6}}$.

\smallskip
\gris{$\Cf$ passe par $(0\,;2)$ ; inégalité large : $0$ est compris.}
\end{lecture}
\begin{figlecture}
\begin{tikzpicture}[x=0.42cm,y=0.42cm]
  \courbesXXXVIII
  \droitey[black!70]{-7}{7}{2}
  \morceau{0}{6}\crochet{0}{[}\crochet{6}{]}
\end{tikzpicture}
\end{figlecture}
\cas{b)}\par\nopagebreak
\begin{lecture}
$g(x)<0$ : $\Cg$ est strictement au-dessous de l'axe des abscisses à droite de $3$ :

$\res{S=\intof{3}{6}}$.

\smallskip
\gris{$g(3)=0$ : exclu ; $g(6)=-1<0$ : inclus.}
\end{lecture}
\begin{figlecture}
\begin{tikzpicture}[x=0.42cm,y=0.42cm]
  \courbesXXXVIII
  \morceau{3}{6}\crochet{3}{]}\crochet{6}{]}
\end{tikzpicture}
\end{figlecture}
\cas{c)}\par\nopagebreak
\begin{lecture}
$g(x)\ge 4$ : $\Cg$ est sur la droite $y=4$ ou au-dessus de $-4$ à $-2$ :

$\res{S=\intff{-4}{-2}}$.

\smallskip
\gris{Le sommet de la bosse est $(-3\,;5)$ ; la courbe passe à la hauteur $4$ en $-4$ et
en $-2$.}
\end{lecture}
\begin{figlecture}
\begin{tikzpicture}[x=0.42cm,y=0.42cm]
  \courbesXXXVIII
  \droitey[black!70]{-7}{7}{4}
  \draw[black!70,dashed] (-4,4)--(-4,0) (-2,4)--(-2,0);
  \morceau{-4}{-2}\crochet{-4}{[}\crochet{-2}{]}
\end{tikzpicture}
\end{figlecture}
\cas{d)}\par\nopagebreak
\begin{lecture}
La droite et la courbe se coupent aux points $(-4\,;4)$, $(2\,;1)$ et $(6\,;-1)$ :

$\res{S=\ens{-4\,;2\,;6}}$.

\smallskip
\gris{Le dernier point commun est l'extrémité droite des deux tracés.}
\end{lecture}
\begin{figlecture}
\begin{tikzpicture}[x=0.42cm,y=0.42cm]
  \courbesXXXVIII
  \path (-4,4) node[croix]{} (2,1) node[croix]{} (6,-1) node[croix]{};
  \draw[black!70,dashed] (-4,4)--(-4,0) (2,1)--(2,0) (6,-1)--(6,0);
  \solx{-4}\solx{2}\solx{6}
\end{tikzpicture}
\end{figlecture}
\cas{e)}\par\nopagebreak
\begin{lecture}
$f(x)<g(x)$ : la droite $\Cf$ est strictement au-dessous de $\Cg$ entre $-4$ et $2$ :

$\res{S=\intoo{-4}{2}}$.

\smallskip
\gris{Inégalité stricte : $-4$ et $2$, où les deux tracés se coupent, sont exclus.}
\end{lecture}
\begin{figlecture}
\begin{tikzpicture}[x=0.42cm,y=0.42cm]
  \courbesXXXVIII
  \draw[black!70,dashed] (-4,4)--(-4,0) (2,1)--(2,0);
  \morceau{-4}{2}\crochet{-4}{]}\crochet{2}{[}
\end{tikzpicture}
\end{figlecture}
\rem{Attention aux couleurs et aux noms : ici, $\Cf$ est la droite et $\Cg$ la courbe.}
\end{exo}

\begin{exo}[Graphique et calcul]
\cas{1)}\par\nopagebreak
\explique{De $-2$ à $3$ avec un pas de $1$ : on calcule chaque image.}{$g(x)=x$ : l'image
d'un nombre est ce nombre lui-même}
\begin{calculs}
f(-2) &= (-2)^{2}-2 = 2 &&\justif{$(-2)^{2}=4$}\\
f(-1) &= (-1)^{2}-2 = -1 &&\\
f(0) &= 0^{2}-2 = -2 &&\\
f(3) &= 3^{2}-2 = 7 &&\justif{de même $f(1)=-1$ et $f(2)=2$}
\end{calculs}
\begin{center}
\renewcommand{\arraystretch}{1.25}
\begin{tabular}{|c|*{6}{c|}}
\hline
$x$ & $-2$ & $-1$ & $0$ & $1$ & $2$ & $3$\\\hline
$f(x)$ & $2$ & $-1$ & $-2$ & $-1$ & $2$ & $7$\\\hline
$g(x)$ & $-2$ & $-1$ & $0$ & $1$ & $2$ & $3$\\\hline
\end{tabular}
\end{center}

\cas{2) et 3)}\par\nopagebreak
\begin{lecture}
2) On place les points des deux tableaux ; ceux de $f$ se relient par une courbe sans
angle, ceux de $g$ sont alignés : $\Cg$ est une droite.

\medskip
3) Les tableaux montrent déjà deux valeurs de $x$ où les images sont égales :
$f(-1)=g(-1)$ et $f(2)=g(2)$. Les courbes se coupent en $(-1\,;-1)$ et $(2\,;2)$ :

$f(x)=g(x)$ : $\res{S=\ens{-1\,;2}}$.

\medskip
$\Cf$ est strictement au-dessous de $\Cg$ entre $-1$ et $2$ :

$f(x)<g(x)$ : $\res{S=\intoo{-1}{2}}$.

\smallskip
\gris{Inégalité stricte : $-1$ et $2$ sont exclus.}
\end{lecture}
\begin{figlecture}
\begin{tikzpicture}[x=0.5cm,y=0.5cm]
  \repere{-3}{-3}{4}{8}
  \gradx{-2,-1,1,2,3}
  \grady{-2,-1,1,2,3,4,5,6,7}
  \draw[coulCourbe,line width=1.1pt,domain=-2:3,samples=60] plot (\x,{(\x)*(\x)-2});
  \draw[coulCourbeC,line width=1.1pt] (-2,-2)--(3,3);
  \foreach \a/\b in {-2/2,-1/-1,0/-2,1/-1,2/2,3/7} \path[coulCourbe] (\a,\b) node[croix]{};
  \foreach \a in {-2,0,1,3} \path[coulCourbeC] (\a,\a) node[croixdroite]{};
  \node[coulCourbe,font=\footnotesize,left] at (2.9,7) {$\Cf$};
  \node[coulCourbeC,font=\footnotesize,right] at (3,3) {$\Cg$};
  \draw[black!70,dashed] (-1,-1)--(-1,0) (2,2)--(2,0);
  \morceau{-1}{2}\crochet{-1}{]}\crochet{2}{[}
\end{tikzpicture}
\end{figlecture}

\cas{4)}\par\nopagebreak
\begin{calculs}
(x+1)(x-2) &= x^{2}-2x+x-2 &&\justif{on développe : $x\times x$ ; $x\times(-2)$ ; $1\times x$ ; $1\times(-2)$}\\
(x+1)(x-2) &= x^{2}-x-2 &&\justif{$-2x+x=-x$ ; c'est bien $x^{2}-2-x$}
\end{calculs}
\begin{calculs}
x^{2}-2 &= x &&\justif{on résout $f(x)=g(x)$}\\
x^{2}-2-x &= 0 &&\justif{on retranche $x$ aux deux membres}\\
(x+1)(x-2) &= 0 &&\justif{la forme factorisée vérifiée au-dessus}
\end{calculs}
\explique{Un produit est nul lorsque l'un de ses facteurs est nul : $x+1=0$ ou $x-2=0$,
c'est-à-dire $x=-1$ ou $x=2$. On retrouve $\res{S=\ens{-1\,;2}}$.}{rappel du collège ; le calcul confirme que
les valeurs lues sont exactes}

\cas{5)}\par\nopagebreak
\begin{calculs}
f(1{,}5) &= 1{,}5^{2}-2 &&\justif{l'image de l'abscisse de $P$}\\
f(1{,}5) &= 2{,}25-2 = 0{,}25 &&
\end{calculs}
\explique{$f(1{,}5)=0{,}25$, l'ordonnée de $P$ : $\res{P\in\Cf}$.}{méthode 12}
\explique{$g(1{,}5)=1{,}5$, alors que l'ordonnée de $P$ est $0{,}25$ :
$\res{P\notin\Cg}$.}{cohérent avec la question 3 : $1{,}5$ est entre $-1$ et $2$, où
$\Cf$ est au-dessous de $\Cg$}
\end{exo}

\end{document}
